%% Generated by Sphinx.
\def\sphinxdocclass{report}
\documentclass[letterpaper,10pt,english]{sphinxmanual}
\ifdefined\pdfpxdimen
   \let\sphinxpxdimen\pdfpxdimen\else\newdimen\sphinxpxdimen
\fi \sphinxpxdimen=.75bp\relax
\ifdefined\pdfimageresolution
    \pdfimageresolution= \numexpr \dimexpr1in\relax/\sphinxpxdimen\relax
\fi
%% let collapsible pdf bookmarks panel have high depth per default
\PassOptionsToPackage{bookmarksdepth=5}{hyperref}

\PassOptionsToPackage{booktabs}{sphinx}
\PassOptionsToPackage{colorrows}{sphinx}

\PassOptionsToPackage{warn}{textcomp}
\usepackage[utf8]{inputenc}
\ifdefined\DeclareUnicodeCharacter
% support both utf8 and utf8x syntaxes
  \ifdefined\DeclareUnicodeCharacterAsOptional
    \def\sphinxDUC#1{\DeclareUnicodeCharacter{"#1}}
  \else
    \let\sphinxDUC\DeclareUnicodeCharacter
  \fi
  \sphinxDUC{00A0}{\nobreakspace}
  \sphinxDUC{2500}{\sphinxunichar{2500}}
  \sphinxDUC{2502}{\sphinxunichar{2502}}
  \sphinxDUC{2514}{\sphinxunichar{2514}}
  \sphinxDUC{251C}{\sphinxunichar{251C}}
  \sphinxDUC{2572}{\textbackslash}
\fi
\usepackage{cmap}
\usepackage[T1]{fontenc}
\usepackage{amsmath,amssymb,amstext}
\usepackage{babel}



\usepackage{tgtermes}
\usepackage{tgheros}
\renewcommand{\ttdefault}{txtt}



\usepackage[Bjarne]{fncychap}
\usepackage{sphinx}

\fvset{fontsize=auto}
\usepackage{geometry}


% Include hyperref last.
\usepackage{hyperref}
% Fix anchor placement for figures with captions.
\usepackage{hypcap}% it must be loaded after hyperref.
% Set up styles of URL: it should be placed after hyperref.
\urlstyle{same}

\addto\captionsenglish{\renewcommand{\contentsname}{Contents}}

\usepackage{sphinxmessages}
\setcounter{tocdepth}{1}



\title{AI\sphinxhyphen{}Based Handover Prediction in UAV\sphinxhyphen{}Assisted 3D Heterogeneous Networks (Beyond 5G / 6G\sphinxhyphen{}Oriented)}
\date{May 31, 2026}
\release{0.1.0}
\author{WASSIM MCHICHI, AOUTMANI MOHAMMED AMIN}
\newcommand{\sphinxlogo}{\vbox{}}
\renewcommand{\releasename}{Release}
\makeindex
\begin{document}

\ifdefined\shorthandoff
  \ifnum\catcode`\=\string=\active\shorthandoff{=}\fi
  \ifnum\catcode`\"=\active\shorthandoff{"}\fi
\fi

\pagestyle{empty}
\sphinxmaketitle
\pagestyle{plain}
\sphinxtableofcontents
\pagestyle{normal}
\phantomsection\label{\detokenize{index::doc}}


\sphinxstepscope


\chapter{Overview: Time\sphinxhyphen{}Series Based Handover Cell Selection for 5G\sphinxhyphen{}to\sphinxhyphen{}6G Migration}
\label{\detokenize{overview:overview-time-series-based-handover-cell-selection-for-5g-to-6g-migration}}\label{\detokenize{overview::doc}}

\section{1. Introduction}
\label{\detokenize{overview:introduction}}
\sphinxAtStartPar
The transition from 5G to 6G requires a paradigm shift in how we manage User Equipment (UE) mobility. As network density and frequency spectrum complexity increase, traditional reactive handover (HO) mechanisms based on instantaneous signal thresholds are becoming obsolete. This study focuses on a \sphinxstylestrong{proactive paradigm}, where handover cell selection is driven by \sphinxstylestrong{time\sphinxhyphen{}series prediction of network indicators}, enabling the network to anticipate movement patterns and channel quality fluctuations before they occur.


\bigskip\hrule\bigskip



\section{2. The Core Problem: Beyond Threshold\sphinxhyphen{}Based Handovers}
\label{\detokenize{overview:the-core-problem-beyond-threshold-based-handovers}}
\sphinxAtStartPar
Current 5G handover protocols rely on the “A3 event” (RSRP\sphinxhyphen{}based triggering). This approach is fundamentally limited in a 6G context for the following reasons:
\begin{itemize}
\item {} 
\sphinxAtStartPar
\sphinxstylestrong{Reactive Latency:} Waiting for signal degradation to cross a threshold creates a “dead zone” of connectivity, incompatible with 6G’s strict URLLC (Ultra\sphinxhyphen{}Reliable Low\sphinxhyphen{}Latency Communication) requirements.

\item {} 
\sphinxAtStartPar
\sphinxstylestrong{Signal Instability:} High\sphinxhyphen{}frequency 6G bands (mmWave/THz) are highly susceptible to blockage and rapid fading. Instantaneous values (current RSRP) are poor predictors of signal viability just 500ms into the future.

\item {} 
\sphinxAtStartPar
\sphinxstylestrong{Oscillation Risks:} Relying on current snapshots leads to frequent “ping\sphinxhyphen{}pong” handovers, as the network fails to distinguish between transient interference and genuine mobility\sphinxhyphen{}driven signal loss.

\end{itemize}


\bigskip\hrule\bigskip



\section{3. The Research Focus: Time\sphinxhyphen{}Series Predictive Decision Making}
\label{\detokenize{overview:the-research-focus-time-series-predictive-decision-making}}
\sphinxAtStartPar
Our research shifts the focus from \sphinxstylestrong{“What is the signal now?”} to \sphinxstylestrong{“What will the indicators be at time \(t + \Delta t\)?”}

\sphinxAtStartPar
By treating network indicators as continuous time\sphinxhyphen{}series data, we can leverage temporal dependencies to make proactive handover decisions.


\subsection{Key Indicators for Prediction:}
\label{\detokenize{overview:key-indicators-for-prediction}}\begin{itemize}
\item {} 
\sphinxAtStartPar
\sphinxstylestrong{Channel State Information (CSI):} Predicting future fading patterns and SNR.

\item {} 
\sphinxAtStartPar
\sphinxstylestrong{Reference Signal Received Power (RSRP):} Forecasting long\sphinxhyphen{}term trends to avoid HOs triggered by momentary shadowing.

\item {} 
\sphinxAtStartPar
\sphinxstylestrong{Network Load/Traffic Density:} Predicting future cell congestion to perform load\sphinxhyphen{}balancing\sphinxhyphen{}aware handovers.

\item {} 
\sphinxAtStartPar
\sphinxstylestrong{UE Trajectory/Velocity:} Using historical coordinate data to estimate the next likely target cells.

\end{itemize}


\bigskip\hrule\bigskip



\section{4. Proposed Framework Architecture}
\label{\detokenize{overview:proposed-framework-architecture}}
\sphinxAtStartPar
The proposed system utilizes a predictive pipeline to inform the selection engine.


\subsection{The Decision\sphinxhyphen{}Making Workflow:}
\label{\detokenize{overview:the-decision-making-workflow}}\begin{enumerate}
\sphinxsetlistlabels{\arabic}{enumi}{enumii}{}{.}%
\item {} 
\sphinxAtStartPar
\sphinxstylestrong{Data Collection:} Continuous streaming of temporal KPI data from the UE and gNB.

\item {} 
\sphinxAtStartPar
\sphinxstylestrong{Temporal Feature Engineering:} Normalizing and windowing data to capture trends, seasonality, and mobility patterns.

\item {} 
\sphinxAtStartPar
\sphinxstylestrong{Predictive Modeling:} Utilizing Deep Learning architectures (e.g., LSTMs, Transformers, or Temporal Convolutional Networks) to forecast indicator values.

\item {} 
\sphinxAtStartPar
\sphinxstylestrong{Proactive Handover Logic:} A decision engine that selects the target cell by comparing the \sphinxstyleemphasis{predicted} states of neighboring cells at time \(t + \Delta t\), rather than their current states.

\end{enumerate}


\bigskip\hrule\bigskip



\section{5. Mathematical Formulation}
\label{\detokenize{overview:mathematical-formulation}}
\sphinxAtStartPar
The objective is to maximize a utility function \(U\) based on predicted indicators \(I\):
\begin{equation*}
\begin{split}
P_{HO} = \arg\max_{c \in C} \mathbb{E}\left[ U\left(\hat{I}_c(t+\Delta t)\right) \right]
\end{split}
\end{equation*}
\sphinxAtStartPar
\sphinxstylestrong{Where:}
\begin{itemize}
\item {} 
\sphinxAtStartPar
\(\hat{I}_c(t+\Delta t)\) is the predicted indicator vector (RSRP, Load, Throughput) for cell \(c\) at the near\sphinxhyphen{}future time step.

\item {} 
\sphinxAtStartPar
\(C\) is the set of candidate neighbor cells.

\item {} 
\sphinxAtStartPar
\(U(\cdot)\) is the multi\sphinxhyphen{}objective utility function, which penalizes:
\begin{itemize}
\item {} 
\sphinxAtStartPar
high handover frequency

\item {} 
\sphinxAtStartPar
low predicted SNR

\end{itemize}

\end{itemize}


\bigskip\hrule\bigskip



\section{6. Strategic Research Goals}
\label{\detokenize{overview:strategic-research-goals}}\begin{itemize}
\item {} 
\sphinxAtStartPar
\sphinxstylestrong{Anticipatory Handover:} Initiate HO preparation \sphinxstyleemphasis{before} the current signal drops below the critical threshold.

\item {} 
\sphinxAtStartPar
\sphinxstylestrong{Ping\sphinxhyphen{}Pong Mitigation:} Filter out noise and transient dips in the time\sphinxhyphen{}series data to reduce redundant HO signaling.

\item {} 
\sphinxAtStartPar
\sphinxstylestrong{Seamless Transition:} Ensure target cell capacity is reserved based on predicted arrival, preventing congestion\sphinxhyphen{}related drops during the handover execution.

\end{itemize}

\sphinxstepscope


\chapter{Dataset Documentation: Geo\sphinxhyphen{}Context Aware Predictive Handover Dataset}
\label{\detokenize{dataset:dataset-documentation-geo-context-aware-predictive-handover-dataset}}\label{\detokenize{dataset::doc}}

\section{1. Executive Summary}
\label{\detokenize{dataset:executive-summary}}
\sphinxAtStartPar
Data generation has been the single most challenging milestone of this project. Evaluating predictive handover cell selection algorithms for 5G NSA\sphinxhyphen{}to\sphinxhyphen{}6G migration requires a highly specific dataset that combines real geographical context, time\sphinxhyphen{}series continuity, dynamic user mobility, and realistic physical layer propagation models capable of handling dual\sphinxhyphen{}connectivity.

\sphinxAtStartPar
To achieve this, we developed a novel framework that mirrors a high\sphinxhyphen{}density urban environment: \sphinxstylestrong{Casablanca Finance City (CFC), Morocco}. The environment leverages real base station coordinates from OpenCellID, maps an E\sphinxhyphen{}UTRA\sphinxhyphen{}NR Dual Connectivity (EN\sphinxhyphen{}DC) architecture, incorporates a custom cell\sphinxhyphen{}load aware dynamic UAV drone control loop, and orchestrates user mobility through a hybrid simulation engine. This document charts our iterative journey across four distinct versions, highlighting the failures, bottlenecks, and the final state\sphinxhyphen{}of\sphinxhyphen{}the\sphinxhyphen{}art hybrid pipeline that made robust multi\sphinxhyphen{}layer time\sphinxhyphen{}series forecasting possible.


\bigskip\hrule\bigskip



\section{2. Core Simulation Entities \& Mechanics}
\label{\detokenize{dataset:core-simulation-entities-mechanics}}

\subsection{2.1. Antenna Densification \& Clustering (Real Geo\sphinxhyphen{}Context)}
\label{\detokenize{dataset:antenna-densification-clustering-real-geo-context}}
\sphinxAtStartPar
To anchor our simulation in reality, we extracted real tower location data from the \sphinxstylestrong{OpenCellID database for Morocco}.
\begin{itemize}
\item {} 
\sphinxAtStartPar
\sphinxstylestrong{Target Area Selection:} Through clustering and densification analysis of national tower deployments, \sphinxstylestrong{Casablanca Finance City (CFC)} was identified as the optimal location representing a high\sphinxhyphen{}density, heterogeneous urban deployment.

\item {} 
\sphinxAtStartPar
\sphinxstylestrong{Pre\sphinxhyphen{}processing:} The raw coordinates were cleaned, filtered for 5G NSA (EN\sphinxhyphen{}DC) deployment topologies, mapping LTE macro\sphinxhyphen{}anchors (Master Nodes) and densified to reflect dense 5G/6G small\sphinxhyphen{}cell overlays (Secondary Nodes).

\end{itemize}


\subsection{2.2. Cell\sphinxhyphen{}Load Aware UAV Drone Control System}
\label{\detokenize{dataset:cell-load-aware-uav-drone-control-system}}
\sphinxAtStartPar
In addition to ground users, our dataset includes simulated Unmanned Aerial Vehicles (UAVs/Drones) behaving as dynamic users/relays. Their movement is not random; it is governed by an automated operational loop:
\begin{itemize}
\item {} 
\sphinxAtStartPar
\sphinxstylestrong{Weighted Centroid Guidance:} Drones track the network’s operational conditions. The target trajectory coordinates are determined by calculating a geometric centroid weighted by the real\sphinxhyphen{}time \sphinxstylestrong{dual\sphinxhyphen{}layer cell load} (combining LTE anchor and 5G active resource utilization) of surrounding ground antennas.

\item {} 
\sphinxAtStartPar
\sphinxstylestrong{Distance Penalization:} The control loop applies a penalty function relative to the distance between the drone and active antennas to avoid extreme, energy\sphinxhyphen{}inefficient trajectories.

\item {} 
\sphinxAtStartPar
\sphinxstylestrong{Control Formulation:}

\end{itemize}
\begin{equation*}
\begin{split}\mathbf{X}_{uav}(t+1) = f\left( \frac{\sum_{i=1}^{N} L_i(t) \cdot \mathbf{X}_i}{\sum_{i=1}^{N} L_i(t)} - \nabla P(\mathbf{D}) \right)\end{split}
\end{equation*}
\sphinxAtStartPar
Where \(L_i(t)\) represents the aggregate dual\sphinxhyphen{}layer cell load of antenna \(i\) (MN/SN) at time \(t\), \(\mathbf{X}_i\) is the antenna position, and \(\nabla P(\mathbf{D})\) is the distance\sphinxhyphen{}based spatial penalization gradient.

\sphinxAtStartPar
\sphinxincludegraphics{{drone_controlle}.png}


\bigskip\hrule\bigskip



\section{3. The 4\sphinxhyphen{}Stage Iterative Dataset Evolution}
\label{\detokenize{dataset:the-4-stage-iterative-dataset-evolution}}
\sphinxAtStartPar
The creation of this dataset required overcoming multiple structural failures. Below is the detailed breakdown of the development lifecycle:

\begin{sphinxVerbatim}[commandchars=\\\{\}]
[v1: 3GPP Toy Script] ──\PYGZgt{} [v2: Script + NS3 Sample] ──\PYGZgt{} [v3: Pure NS3 Engine] ──\PYGZgt{} [v4: Hybrid Production Pipeline]
     (Too Simple)              (Stats Test Failed)         (Imbalance/CPU Bound)         (Hexagonal + Ray\PYGZhy{}Tracing)
\end{sphinxVerbatim}


\subsection{Version 1: 3GPP Hardcoded Script (Toy Dataset)}
\label{\detokenize{dataset:version-1-3gpp-hardcoded-script-toy-dataset}}\begin{itemize}
\item {} 
\sphinxAtStartPar
\sphinxstylestrong{Methodology:} A standalone Python script designed to output KPI metrics based directly on closed\sphinxhyphen{}form 3GPP standard propagation equations for LTE and 5G bands.

\item {} 
\sphinxAtStartPar
\sphinxstylestrong{Bottlenecks \& Failures:} The resulting data was entirely deterministic and low\sphinxhyphen{}dimensional. It completely lacked the stochastic variance, multi\sphinxhyphen{}path fading, interlocking dual\sphinxhyphen{}connectivity dependencies, and complex user mobility profiles of real\sphinxhyphen{}world networks. It acted as a “toy dataset” and was useless for deep learning.

\end{itemize}


\subsection{Version 2: Hybrid Script + NS3 Sample (Statistical Validation Failure)}
\label{\detokenize{dataset:version-2-hybrid-script-ns3-sample-statistical-validation-failure}}\begin{itemize}
\item {} 
\sphinxAtStartPar
\sphinxstylestrong{Methodology:} In an attempt to add realism, we generated a small benchmark sample using the NS\sphinxhyphen{}3 (Network Simulator 3) EN\sphinxhyphen{}DC network stack and augmented it using our internal Python hardcoded script. We applied rigorous statistical integrity tests (e.g., distribution matching, correlation alignment) to check if the script could scale up the dataset.

\item {} 
\sphinxAtStartPar
\sphinxstylestrong{Bottlenecks \& Failures:} \sphinxstylestrong{The majority of statistical tests failed.} The Python script was mathematically too simplistic to mimic the intricate, stateful behavior of the NS\sphinxhyphen{}3 discrete\sphinxhyphen{}event network simulator (such as coupled LTE/5G MAC\sphinxhyphen{}layer scheduling overhead, EN\sphinxhyphen{}DC bearer splits, and protocol stack delays).

\end{itemize}


\subsection{Version 3: Pure NS\sphinxhyphen{}3 Dataset Generation (Imbalance \& Compute Bottlenecks)}
\label{\detokenize{dataset:version-3-pure-ns-3-dataset-generation-imbalance-compute-bottlenecks}}\begin{itemize}
\item {} 
\sphinxAtStartPar
\sphinxstylestrong{Methodology:} We abandoned the hardcoded scripts and moved the entire simulation infrastructure directly inside \sphinxstylestrong{NS\sphinxhyphen{}3}.

\item {} 
\sphinxAtStartPar
\sphinxstylestrong{Bottlenecks \& Critical Failures:} This version introduced two catastrophic blockers:

\end{itemize}
\begin{enumerate}
\sphinxsetlistlabels{\arabic}{enumi}{enumii}{}{.}%
\item {} 
\sphinxAtStartPar
NS\sphinxhyphen{}3 System\sphinxhyphen{}Level Abstraction \& Artificial Event Sparsity

\end{enumerate}

\sphinxAtStartPar
As a system\sphinxhyphen{}level simulator, NS\sphinxhyphen{}3 fundamentally relies on abstracted, distance\sphinxhyphen{}smoothed statistical propagation models rather than environment\sphinxhyphen{}specific physical layer dynamics. Consequently, its simulated RF environment degrades predictably and smoothly. Because the simulation lacks authentic Line\sphinxhyphen{}of\sphinxhyphen{}Sight (LOS) to Non\sphinxhyphen{}Line\sphinxhyphen{}of\sphinxhyphen{}Sight (NLOS) transitions and sudden urban blockage effects, the signal rarely experienced the steep, realistic drops that trigger dual\sphinxhyphen{}connectivity handovers (MN/SN changes) in the real world.

\sphinxAtStartPar
This architectural limitation resulted in an artificially low handover rate (e.g., \textasciitilde{}0.7\%), meaning handovers only occurred at extreme geometric cell edges. This simulator\sphinxhyphen{}induced sparsity created a severe class imbalance in our dataset, rendering traditional machine learning fixes useless:
\begin{itemize}
\item {} 
\sphinxAtStartPar
Focal Loss: Proved mathematically insufficient to overcome the severe classification bias during gradient descent.

\item {} 
\sphinxAtStartPar
SMOTE (Synthetic Minority Over\sphinxhyphen{}sampling Technique): Proved fundamentally incompatible with our architecture. SMOTE interpolates between isolated feature vectors, which destroys the critical temporal continuity and sequential phase alignment of time\sphinxhyphen{}series data. It injected random noise into our signals, rendering sequential models (LSTMs/Transformers) unable to learn true pre\sphinxhyphen{}handover degradation signatures.

\end{itemize}
\begin{enumerate}
\sphinxsetlistlabels{\arabic}{enumi}{enumii}{}{.}%
\setcounter{enumi}{1}
\item {} 
\sphinxAtStartPar
\sphinxstylestrong{Computational Complexity (“Time is Gold”):} NS\sphinxhyphen{}3 is purely CPU\sphinxhyphen{}bound and single\sphinxhyphen{}threaded by nature for synchronous event execution. Modeling concurrent LTE and 5G stacks meant a single simulation run took \sphinxstylestrong{at least 8 hours}, killing our iteration speed.

\end{enumerate}


\subsection{Version 4: The Production Pipeline (Hybrid Framework)}
\label{\detokenize{dataset:version-4-the-production-pipeline-hybrid-framework}}\begin{itemize}
\item {} 
\sphinxAtStartPar
\sphinxstylestrong{Methodology:} Our final breakthrough utilizes a decoupled, multi\sphinxhyphen{}engine production pipeline that splits mobility, control loops, and high\sphinxhyphen{}fidelity physical layer calculations across specialized tools.

\end{itemize}


\begin{savenotes}\sphinxattablestart
\sphinxthistablewithglobalstyle
\centering
\begin{tabulary}{\linewidth}[t]{TTT}
\sphinxtoprule
\sphinxstyletheadfamily 
\sphinxAtStartPar
Layer
&\sphinxstyletheadfamily 
\sphinxAtStartPar
Orchestrating Technology
&\sphinxstyletheadfamily 
\sphinxAtStartPar
Functional Responsibility
\\
\sphinxmidrule
\sphinxtableatstartofbodyhook
\sphinxAtStartPar
\sphinxstylestrong{Mobility}
&
\sphinxAtStartPar
NS\sphinxhyphen{}3 (Network Simulator 3)
&
\sphinxAtStartPar
Coordinates ground user trajectories and spatial transitions within the EN\sphinxhyphen{}DC footprint.
\\
\sphinxhline
\sphinxAtStartPar
\sphinxstylestrong{Control}
&
\sphinxAtStartPar
Custom Python Engine
&
\sphinxAtStartPar
Manages the cell\sphinxhyphen{}load aware dynamic UAV drone navigation loops.
\\
\sphinxhline
\sphinxAtStartPar
\sphinxstylestrong{RF / Signal}
&
\sphinxAtStartPar
Sionna (GPU\sphinxhyphen{}Accelerated)
&
\sphinxAtStartPar
Executes massive parallel physical layer, fading, and noise calculations across heterogeneous frequencies (LTE and 5G).
\\
\sphinxbottomrule
\end{tabulary}
\sphinxtableafterendhook\par
\sphinxattableend\end{savenotes}

\sphinxAtStartPar
To manage scale, compute costs, and model accuracy, Version 4 is divided into \sphinxstylestrong{two structural sub\sphinxhyphen{}datasets}:


\subsubsection{Sub\sphinxhyphen{}Dataset A: \sphinxstyleliteralintitle{\sphinxupquote{hexagonal\_simulation}}}
\label{\detokenize{dataset:sub-dataset-a-hexagonal-simulation}}\begin{itemize}
\item {} 
\sphinxAtStartPar
\sphinxstylestrong{Characteristics:} High\sphinxhyphen{}volume, statistical channel modeling.

\item {} 
\sphinxAtStartPar
\sphinxstylestrong{Implementation:} Built using Sionna’s \sphinxstylestrong{UMIChannel (Urban Micro)} model. It incorporates comprehensive multi\sphinxhyphen{}cell co\sphinxhyphen{}channel interference modeling for both LTE and 5G layers, fast/slow fading profiles, log\sphinxhyphen{}normal shadowing, and thermal noise layers.

\item {} 
\sphinxAtStartPar
\sphinxstylestrong{Purpose:} This dataset serves as the backbone for baseline training, teaching the SOTA neural network architectures generalized macro\sphinxhyphen{}temporal dependencies and basic multi\sphinxhyphen{}connectivity handover behavior.

\end{itemize}


\subsubsection{Sub\sphinxhyphen{}Dataset B: \sphinxstyleliteralintitle{\sphinxupquote{ray\_tracing\_simulation}}}
\label{\detokenize{dataset:sub-dataset-b-ray-tracing-simulation}}\begin{itemize}
\item {} 
\sphinxAtStartPar
\sphinxstylestrong{Characteristics:} High\sphinxhyphen{}fidelity, deterministic, environment\sphinxhyphen{}specific modeling.

\item {} 
\sphinxAtStartPar
\sphinxstylestrong{Implementation:} We imported a highly detailed \sphinxstylestrong{OpenStreetMap (\sphinxcode{\sphinxupquote{.osm}}) file of Casablanca Finance City (CFC)} into \sphinxstylestrong{Blender}. We manually configured realistic 3D building heights and accurately distributed the three most common urban construction materials across the scene.

\item {} 
\sphinxAtStartPar
\sphinxstylestrong{The Ray\sphinxhyphen{}Tracing Pipeline:} This 3D environment was fed into Sionna’s GPU\sphinxhyphen{}accelerated Ray Tracing engine to simulate real physical interactions (diffraction, specular reflection, and scattering).

\item {} 
\sphinxAtStartPar
\sphinxstylestrong{Purpose (Fine\sphinxhyphen{}Tuning Paradigm):} Because ray tracing is computationally expensive, we generate data in small, highly controlled batches. The base model, pre\sphinxhyphen{}trained on the large\sphinxhyphen{}scale \sphinxcode{\sphinxupquote{hexagonal\_simulation}}, uses this deterministic ray\sphinxhyphen{}tracing dataset for \sphinxstylestrong{fine\sphinxhyphen{}tuning}, allowing it to adapt to real\sphinxhyphen{}world geometric multipath anomalies and localized building shadowing profiles.

\end{itemize}


\begin{savenotes}\sphinxattablestart
\sphinxthistablewithglobalstyle
\centering
\begin{tabulary}{\linewidth}[t]{TT}
\sphinxtoprule
\sphinxstyletheadfamily 
\sphinxAtStartPar
Image 1
&\sphinxstyletheadfamily 
\sphinxAtStartPar
Image 2
\\
\sphinxmidrule
\sphinxtableatstartofbodyhook
\sphinxAtStartPar
\sphinxincludegraphics{{city_raytracing_blender}.png}
&
\sphinxAtStartPar
\sphinxincludegraphics{{drone_ray_tracin_blender}.png}
\\
\sphinxbottomrule
\end{tabulary}
\sphinxtableafterendhook\par
\sphinxattableend\end{savenotes}


\begin{savenotes}\sphinxattablestart
\sphinxthistablewithglobalstyle
\centering
\begin{tabulary}{\linewidth}[t]{TT}
\sphinxtoprule
\sphinxstyletheadfamily 
\sphinxAtStartPar
Image 3
&\sphinxstyletheadfamily 
\sphinxAtStartPar
Image 4
\\
\sphinxmidrule
\sphinxtableatstartofbodyhook
\sphinxAtStartPar
\sphinxincludegraphics{{render-ray-tracing-0001}.png}
&
\sphinxAtStartPar
\sphinxincludegraphics{{render-ray-tracing-0002}.png}
\\
\sphinxbottomrule
\end{tabulary}
\sphinxtableafterendhook\par
\sphinxattableend\end{savenotes}


\begin{savenotes}\sphinxattablestart
\sphinxthistablewithglobalstyle
\centering
\begin{tabulary}{\linewidth}[t]{TT}
\sphinxtoprule
\sphinxstyletheadfamily 
\sphinxAtStartPar
Image 5
&\sphinxstyletheadfamily 
\sphinxAtStartPar
Image 6
\\
\sphinxmidrule
\sphinxtableatstartofbodyhook
\sphinxAtStartPar
\sphinxincludegraphics{{render-ray-tracing-0003}.png}
&
\sphinxAtStartPar
\sphinxincludegraphics{{user_ray_tracing_blender}.png}
\\
\sphinxbottomrule
\end{tabulary}
\sphinxtableafterendhook\par
\sphinxattableend\end{savenotes}


\bigskip\hrule\bigskip



\section{4. Time\sphinxhyphen{}Series Feature Schema}
\label{\detokenize{dataset:time-series-feature-schema}}
\sphinxAtStartPar
The final dataset is structured as sliding windows of sequence length \(T\), collecting the following key features:

\begin{sphinxVerbatim}[commandchars=\\\{\}]
[ t\PYGZhy{}(T\PYGZhy{}1) , ... , t\PYGZhy{}1 , t ]  ───\PYGZgt{} Predict Target Event at ───\PYGZgt{} [ t + Δt ]
\end{sphinxVerbatim}


\subsection{4.1. Feature Matrix Specification}
\label{\detokenize{dataset:feature-matrix-specification}}\begin{enumerate}
\sphinxsetlistlabels{\arabic}{enumi}{enumii}{}{.}%
\item {} 
\sphinxAtStartPar
\sphinxstylestrong{Temporal RF Metrics:} Time\sphinxhyphen{}series arrays of RSRP, RSRQ, and SINR for both the serving LTE Master Node (MN) and 5G Secondary Node (SN), plus top \(N\) neighboring candidate cells for each layer (\(T \times 2(N+1)\) matrix dimensions).

\item {} 
\sphinxAtStartPar
\sphinxstylestrong{Network State Indicators:} Real\sphinxhyphen{}time cell load, queue delay profiles, and localized throughput capacities per base station (spanning both LTE anchors and 5G nodes).

\item {} 
\sphinxAtStartPar
\sphinxstylestrong{Mobility Metrics:} Multi\sphinxhyphen{}axis velocity vectors (\(v_x, v_y, v_z\)) and spatial trajectories of ground users and dynamic UAV centroids.

\item {} 
\sphinxAtStartPar
\sphinxstylestrong{Target Label:} Explicit handover event vectors marked at time \(t + \Delta t\), identifying the optimal target action (e.g., MN Handover or SN Change) while avoiding ping\sphinxhyphen{}pong oscillations.

\end{enumerate}

\sphinxstepscope


\chapter{Model Architecture Reference}
\label{\detokenize{architecture:model-architecture-reference}}\label{\detokenize{architecture::doc}}
\sphinxAtStartPar
This page documents every model architecture developed in the experimentation pipeline, from the first baseline proof\sphinxhyphen{}of\sphinxhyphen{}concept through to the current \sphinxstylestrong{SOTA production model}. All architectures operate on the same core problem: given a sliding window of RF measurements across a candidate cell set, predict the optimal target cell for a handover 5 steps into the future.
\begin{quote}

\sphinxAtStartPar
\sphinxstylestrong{Production model:} \sphinxcode{\sphinxupquote{models/best\_ray\_id\_crossformer.keras}}\\
\sphinxstylestrong{Problem type (production):} 5\sphinxhyphen{}step future horizon multi\sphinxhyphen{}class cell selection\\
\sphinxstylestrong{Problem type (baseline only):} Single\sphinxhyphen{}step (current\sphinxhyphen{}instant) selection — \sphinxcode{\sphinxupquote{notebooks/modeling/01\_temporal\_deepset.ipynb}} only
\end{quote}


\bigskip\hrule\bigskip



\section{Architecture Lineage Overview}
\label{\detokenize{architecture:architecture-lineage-overview}}
\begin{sphinxVerbatim}[commandchars=\\\{\}]
Exp 01 — Temporal DeepSet (baseline, single\PYGZhy{}step)
      │
      ├─► Exp 02 — SpatioTemporal DeepSet Production (5\PYGZhy{}step, multi\PYGZhy{}horizon)
      │
      ├─► Exp 03 — MTL Set Transformer (multi\PYGZhy{}task, 5\PYGZhy{}step)
      │
      ├─► Exp 04 — 6G Predictive Set Transformer (Optuna\PYGZhy{}tuned, 5\PYGZhy{}step)
      │
      ├─► Exp 05 — Strategic DeepSet (context injection, 5\PYGZhy{}step)
      │
      ├─► Exp 06 — SOTA Ray\PYGZhy{}Tracing Fine\PYGZhy{}Tune (MTL backbone \(\rightarrow\) ray domain)
      │
      └─► Exp 07 — Ray\PYGZhy{}ID CrossFormer ★ SOTA / Production (5\PYGZhy{}step)
\end{sphinxVerbatim}


\bigskip\hrule\bigskip



\section{Exp 01 — Temporal DeepSet (Baseline)}
\label{\detokenize{architecture:exp-01-temporal-deepset-baseline}}\begin{quote}

\sphinxAtStartPar
\sphinxstylestrong{Notebook:} \sphinxcode{\sphinxupquote{notebooks/modeling/01\_temporal\_deepset.ipynb}}\\
\sphinxstylestrong{Checkpoint:} \sphinxcode{\sphinxupquote{models/best\_temporal\_deepset.keras}}\\
\sphinxstylestrong{⚠️ Important:} This is a \sphinxstylestrong{single\sphinxhyphen{}step (current\sphinxhyphen{}instant) cell selection} experiment — it predicts the best cell at time \sphinxstyleemphasis{t}, not a future horizon. All subsequent experiments (02\textendash{}07) target 5\sphinxhyphen{}step future prediction for production use.
\end{quote}


\subsection{Architecture}
\label{\detokenize{architecture:architecture}}
\sphinxAtStartPar
The Temporal DeepSet stacks a per\sphinxhyphen{}cell LSTM temporal encoder with a DeepSet permutation\sphinxhyphen{}invariant aggregation, making it robust to variable numbers of neighbor cells and to the arbitrary ordering of the candidate set.

\begin{sphinxVerbatim}[commandchars=\\\{\}]
Input: (K, T, F)  — K cells, T=25 timesteps, F=3 features (RSRP, SINR, Load)
    │
    ├─ [Per\PYGZhy{}Cell LSTM]  (shared weights across K)  \(\rightarrow\)  φ(xᵢ) ∈ ℝ\PYGZca{}64
    │
    ├─ [MaskedGlobalAveragePooling]   \(\rightarrow\)  ρ\PYGZhy{}aggregation over valid cells
    │
    └─ [Classifier MLP]  \(\rightarrow\)  logits ∈ ℝ\PYGZca{}K
\end{sphinxVerbatim}

\sphinxAtStartPar
\sphinxincludegraphics{{temporal_deepset_architecture}.png}


\subsection{Hyperparameters}
\label{\detokenize{architecture:hyperparameters}}

\begin{savenotes}\sphinxattablestart
\sphinxthistablewithglobalstyle
\centering
\begin{tabulary}{\linewidth}[t]{TT}
\sphinxtoprule
\sphinxstyletheadfamily 
\sphinxAtStartPar
Parameter
&\sphinxstyletheadfamily 
\sphinxAtStartPar
Value
\\
\sphinxmidrule
\sphinxtableatstartofbodyhook
\sphinxAtStartPar
\sphinxcode{\sphinxupquote{OBS\_STEPS}}
&
\sphinxAtStartPar
25
\\
\sphinxhline
\sphinxAtStartPar
\sphinxcode{\sphinxupquote{MAX\_CELLS}}
&
\sphinxAtStartPar
10
\\
\sphinxhline
\sphinxAtStartPar
\sphinxcode{\sphinxupquote{LSTM\_UNITS}}
&
\sphinxAtStartPar
64
\\
\sphinxhline
\sphinxAtStartPar
\sphinxcode{\sphinxupquote{PHI\_DIM}}
&
\sphinxAtStartPar
64
\\
\sphinxhline
\sphinxAtStartPar
\sphinxcode{\sphinxupquote{PHI\_LAYERS}}
&
\sphinxAtStartPar
2
\\
\sphinxhline
\sphinxAtStartPar
\sphinxcode{\sphinxupquote{DROPOUT}}
&
\sphinxAtStartPar
0.25
\\
\sphinxhline
\sphinxAtStartPar
\sphinxcode{\sphinxupquote{BATCH\_SIZE}}
&
\sphinxAtStartPar
128
\\
\sphinxhline
\sphinxAtStartPar
\sphinxcode{\sphinxupquote{LR\_INIT}}
&
\sphinxAtStartPar
1e\sphinxhyphen{}3
\\
\sphinxhline
\sphinxAtStartPar
Loss
&
\sphinxAtStartPar
Focal (γ=2.0, α=0.25) + Label Smoothing 0.1
\\
\sphinxbottomrule
\end{tabulary}
\sphinxtableafterendhook\par
\sphinxattableend\end{savenotes}


\subsection{Results (Single\sphinxhyphen{}Step — Baseline Only)}
\label{\detokenize{architecture:results-single-step-baseline-only}}

\begin{savenotes}\sphinxattablestart
\sphinxthistablewithglobalstyle
\centering
\begin{tabulary}{\linewidth}[t]{TT}
\sphinxtoprule
\sphinxstyletheadfamily 
\sphinxAtStartPar
Metric
&\sphinxstyletheadfamily 
\sphinxAtStartPar
Value
\\
\sphinxmidrule
\sphinxtableatstartofbodyhook
\sphinxAtStartPar
\sphinxstylestrong{Test Top\sphinxhyphen{}1}
&
\sphinxAtStartPar
\sphinxstylestrong{99.82 \%}
\\
\sphinxhline
\sphinxAtStartPar
\sphinxstylestrong{Test Top\sphinxhyphen{}3}
&
\sphinxAtStartPar
\sphinxstylestrong{100.0 \%}
\\
\sphinxhline
\sphinxAtStartPar
\sphinxstylestrong{Test Top\sphinxhyphen{}5}
&
\sphinxAtStartPar
\sphinxstylestrong{100.0 \%}
\\
\sphinxbottomrule
\end{tabulary}
\sphinxtableafterendhook\par
\sphinxattableend\end{savenotes}
\begin{quote}

\sphinxAtStartPar
These near\sphinxhyphen{}perfect numbers reflect the \sphinxstylestrong{single\sphinxhyphen{}step} nature of this experiment: the model selects the best \sphinxstyleemphasis{current} cell, which is an easier problem. This is \sphinxstylestrong{not} the 5\sphinxhyphen{}step production task.
\end{quote}


\subsection{Training Curves}
\label{\detokenize{architecture:training-curves}}
\sphinxAtStartPar
\sphinxincludegraphics{{temporal_deepset_training_curves}.png}


\subsection{Confusion Matrix}
\label{\detokenize{architecture:confusion-matrix}}
\sphinxAtStartPar
\sphinxincludegraphics{{temporal_deepset_confusion_matrix}.png}


\bigskip\hrule\bigskip



\section{Exp 02 — SpatioTemporal DeepSet Production}
\label{\detokenize{architecture:exp-02-spatiotemporal-deepset-production}}\begin{quote}

\sphinxAtStartPar
\sphinxstylestrong{Notebook:} \sphinxcode{\sphinxupquote{notebooks/modeling/02\_temporal\_deepset\_production.ipynb}}\\
\sphinxstylestrong{Architecture class:} \sphinxcode{\sphinxupquote{SpatioTemporalDeepSet}}\\
\sphinxstylestrong{Checkpoint:} \sphinxcode{\sphinxupquote{models/best\_temporal\_deepset.keras}}
\end{quote}

\sphinxAtStartPar
This version extends the baseline to the 5\sphinxhyphen{}step future\sphinxhyphen{}horizon prediction problem. A multi\sphinxhyphen{}task decoder simultaneously predicts:
\begin{itemize}
\item {} 
\sphinxAtStartPar
the optimal cell at the \sphinxstylestrong{current instant} (auxiliary task, easy),

\item {} 
\sphinxAtStartPar
the optimal cell at \sphinxstylestrong{each of the next 5 steps} (primary task),

\item {} 
\sphinxAtStartPar
future \sphinxstylestrong{RSRP / SINR / Load signal regression} (auxiliary regression head).

\end{itemize}

\begin{sphinxVerbatim}[commandchars=\\\{\}]
Input: (K, T=200, F=5)
    │
    ├─ [Per\PYGZhy{}Cell LSTM (128 units)] \(\rightarrow\) cell embedding φᵢ ∈ ℝ\PYGZca{}64
    │
    ├─ [Attend module] \(\rightarrow\) context\PYGZhy{}aware aggregation ρ ∈ ℝ\PYGZca{}128
    │
    ├─ [Now\PYGZhy{}Decoder MLP] \(\rightarrow\) cell\PYGZus{}now logits ∈ ℝ\PYGZca{}K         (auxiliary)
    │
    ├─ [Future\PYGZhy{}Decoder LSTM (256)] \(\rightarrow\) hidden per horizon
    │       └─ [Scorer] \(\rightarrow\) logits ∈ ℝ\PYGZca{}(H×K)              (primary, H=5)
    │
    └─ [Signal Regression Head] \(\rightarrow\) (RSRP, SINR, Load) × H  (auxiliary)
\end{sphinxVerbatim}


\subsection{Results (5\sphinxhyphen{}Step Future Prediction)}
\label{\detokenize{architecture:results-5-step-future-prediction}}

\begin{savenotes}\sphinxattablestart
\sphinxthistablewithglobalstyle
\centering
\begin{tabulary}{\linewidth}[t]{TT}
\sphinxtoprule
\sphinxstyletheadfamily 
\sphinxAtStartPar
Horizon
&\sphinxstyletheadfamily 
\sphinxAtStartPar
Top\sphinxhyphen{}1
\\
\sphinxmidrule
\sphinxtableatstartofbodyhook
\sphinxAtStartPar
t+1
&
\sphinxAtStartPar
64.98 \%
\\
\sphinxhline
\sphinxAtStartPar
t+2
&
\sphinxAtStartPar
57.94 \%
\\
\sphinxhline
\sphinxAtStartPar
t+3
&
\sphinxAtStartPar
54.26 \%
\\
\sphinxhline
\sphinxAtStartPar
t+4
&
\sphinxAtStartPar
52.29 \%
\\
\sphinxhline
\sphinxAtStartPar
\sphinxstylestrong{t+5}
&
\sphinxAtStartPar
\sphinxstylestrong{51.20 \%}
\\
\sphinxhline
\sphinxAtStartPar
\sphinxstylestrong{avg}
&
\sphinxAtStartPar
\sphinxstylestrong{56.13 \%}
\\
\sphinxbottomrule
\end{tabulary}
\sphinxtableafterendhook\par
\sphinxattableend\end{savenotes}


\begin{savenotes}\sphinxattablestart
\sphinxthistablewithglobalstyle
\centering
\begin{tabulary}{\linewidth}[t]{TT}
\sphinxtoprule
\sphinxstyletheadfamily 
\sphinxAtStartPar
Metric
&\sphinxstyletheadfamily 
\sphinxAtStartPar
Value
\\
\sphinxmidrule
\sphinxtableatstartofbodyhook
\sphinxAtStartPar
RSRP MAE
&
\sphinxAtStartPar
5.56 dBm
\\
\sphinxhline
\sphinxAtStartPar
SINR MAE
&
\sphinxAtStartPar
5.14 dB
\\
\sphinxhline
\sphinxAtStartPar
Load MAE
&
\sphinxAtStartPar
0.072
\\
\sphinxbottomrule
\end{tabulary}
\sphinxtableafterendhook\par
\sphinxattableend\end{savenotes}


\subsection{Training Curves}
\label{\detokenize{architecture:id1}}
\sphinxAtStartPar
\sphinxincludegraphics{{production_deepset_training_curves}.png}


\subsection{Per\sphinxhyphen{}Step Future Accuracy}
\label{\detokenize{architecture:per-step-future-accuracy}}
\sphinxAtStartPar
\sphinxincludegraphics{{future_step_accuracy}.png}


\bigskip\hrule\bigskip



\section{Exp 03 — MTL Set Transformer}
\label{\detokenize{architecture:exp-03-mtl-set-transformer}}\begin{quote}

\sphinxAtStartPar
\sphinxstylestrong{Notebook:} \sphinxcode{\sphinxupquote{notebooks/modeling/03\_mtl\_transformer.ipynb}}\\
\sphinxstylestrong{Checkpoint:} \sphinxcode{\sphinxupquote{models/best\_mtl\_transformer.keras}}\\
\sphinxstylestrong{Loss:} \sphinxcode{\sphinxupquote{1.0 × Focal(γ=2.0, α=0.25) + 0.5 × MSE}}
\end{quote}

\sphinxAtStartPar
The MTL Set Transformer replaces the DeepSet ρ\sphinxhyphen{}aggregation with a full multi\sphinxhyphen{}head self\sphinxhyphen{}attention block over the cell axis. This enables richer cross\sphinxhyphen{}cell interaction modeling before the final classification head.

\sphinxAtStartPar
\sphinxincludegraphics{{mtl_transformer_architecture}.png}

\begin{sphinxVerbatim}[commandchars=\\\{\}]
Input: (K, T=25, F=4)   [RSRP, SINR, Load + zero\PYGZhy{}anchor flag]
    │
    ├─ [Per\PYGZhy{}Cell LSTM (64)] \(\rightarrow\) φᵢ ∈ ℝ\PYGZca{}64
    │
    ├─ [Set Transformer Block × 2]
    │       ├─ Multi\PYGZhy{}Head Self\PYGZhy{}Attention (4 heads, key\PYGZus{}dim=16)
    │       └─ Feed\PYGZhy{}Forward (FF\PYGZus{}DIM=128)
    │
    ├─ [Cell Classifier Head] \(\rightarrow\) logits ∈ ℝ\PYGZca{}K              (5\PYGZhy{}step)
    │
    └─ [RSRP Regression Head] \(\rightarrow\) predicted RSRP per horizon
\end{sphinxVerbatim}


\subsection{Results}
\label{\detokenize{architecture:results}}

\begin{savenotes}\sphinxattablestart
\sphinxthistablewithglobalstyle
\centering
\begin{tabulary}{\linewidth}[t]{TT}
\sphinxtoprule
\sphinxstyletheadfamily 
\sphinxAtStartPar
Metric
&\sphinxstyletheadfamily 
\sphinxAtStartPar
Value
\\
\sphinxmidrule
\sphinxtableatstartofbodyhook
\sphinxAtStartPar
\sphinxstylestrong{Test Top\sphinxhyphen{}1}
&
\sphinxAtStartPar
\sphinxstylestrong{55.79 \%}
\\
\sphinxhline
\sphinxAtStartPar
\sphinxstylestrong{Test Top\sphinxhyphen{}3}
&
\sphinxAtStartPar
\sphinxstylestrong{86.02 \%}
\\
\sphinxhline
\sphinxAtStartPar
\sphinxstylestrong{Test Top\sphinxhyphen{}5}
&
\sphinxAtStartPar
\sphinxstylestrong{94.99 \%}
\\
\sphinxhline
\sphinxAtStartPar
RSRP MAE
&
\sphinxAtStartPar
4.20 dBm
\\
\sphinxbottomrule
\end{tabulary}
\sphinxtableafterendhook\par
\sphinxattableend\end{savenotes}


\subsection{Training Curves}
\label{\detokenize{architecture:id2}}
\sphinxAtStartPar
\sphinxincludegraphics{{mtl_training_curves}.png}

\sphinxAtStartPar
\sphinxincludegraphics{{mtl_exp3_results}.png}


\bigskip\hrule\bigskip



\section{Exp 04 — 6G Predictive Set Transformer}
\label{\detokenize{architecture:exp-04-6g-predictive-set-transformer}}\begin{quote}

\sphinxAtStartPar
\sphinxstylestrong{Notebook:} \sphinxcode{\sphinxupquote{notebooks/modeling/04\_6g\_predictive.ipynb}}\\
\sphinxstylestrong{Checkpoint:} \sphinxcode{\sphinxupquote{models/best\_6g\_predictive.keras}}\\
\sphinxstylestrong{Tuning:} Optuna hyperparameter search
\end{quote}

\sphinxAtStartPar
This experiment applies automated Optuna\sphinxhyphen{}based hyperparameter optimization over the Set Transformer family, extending the observation window to 200 timesteps to capture long\sphinxhyphen{}range temporal dependencies.

\sphinxAtStartPar
\sphinxincludegraphics{{set_transformer_architecture}.png}


\subsection{Hyperparameters (Optuna Best)}
\label{\detokenize{architecture:hyperparameters-optuna-best}}

\begin{savenotes}\sphinxattablestart
\sphinxthistablewithglobalstyle
\centering
\begin{tabulary}{\linewidth}[t]{TT}
\sphinxtoprule
\sphinxstyletheadfamily 
\sphinxAtStartPar
Parameter
&\sphinxstyletheadfamily 
\sphinxAtStartPar
Value
\\
\sphinxmidrule
\sphinxtableatstartofbodyhook
\sphinxAtStartPar
\sphinxcode{\sphinxupquote{OBS\_STEPS}}
&
\sphinxAtStartPar
200
\\
\sphinxhline
\sphinxAtStartPar
\sphinxcode{\sphinxupquote{LSTM\_UNITS}}
&
\sphinxAtStartPar
256
\\
\sphinxhline
\sphinxAtStartPar
\sphinxcode{\sphinxupquote{N\_HEADS}}
&
\sphinxAtStartPar
2
\\
\sphinxhline
\sphinxAtStartPar
\sphinxcode{\sphinxupquote{N\_ST\_BLOCKS}}
&
\sphinxAtStartPar
2
\\
\sphinxhline
\sphinxAtStartPar
\sphinxcode{\sphinxupquote{FF\_DIM}}
&
\sphinxAtStartPar
128
\\
\sphinxhline
\sphinxAtStartPar
\sphinxcode{\sphinxupquote{DROPOUT}}
&
\sphinxAtStartPar
0.378
\\
\sphinxhline
\sphinxAtStartPar
\sphinxcode{\sphinxupquote{LR\_INIT}}
&
\sphinxAtStartPar
1.68e\sphinxhyphen{}3
\\
\sphinxhline
\sphinxAtStartPar
\sphinxcode{\sphinxupquote{LAMBDA\_CLS}}
&
\sphinxAtStartPar
1.0
\\
\sphinxhline
\sphinxAtStartPar
\sphinxcode{\sphinxupquote{LAMBDA\_REG}}
&
\sphinxAtStartPar
0.7
\\
\sphinxbottomrule
\end{tabulary}
\sphinxtableafterendhook\par
\sphinxattableend\end{savenotes}


\subsection{Results}
\label{\detokenize{architecture:id3}}

\begin{savenotes}\sphinxattablestart
\sphinxthistablewithglobalstyle
\centering
\begin{tabulary}{\linewidth}[t]{TT}
\sphinxtoprule
\sphinxstyletheadfamily 
\sphinxAtStartPar
Metric
&\sphinxstyletheadfamily 
\sphinxAtStartPar
Value
\\
\sphinxmidrule
\sphinxtableatstartofbodyhook
\sphinxAtStartPar
\sphinxstylestrong{Test Top\sphinxhyphen{}1}
&
\sphinxAtStartPar
\sphinxstylestrong{54.93 \%}
\\
\sphinxhline
\sphinxAtStartPar
\sphinxstylestrong{Test Top\sphinxhyphen{}3}
&
\sphinxAtStartPar
\sphinxstylestrong{85.47 \%}
\\
\sphinxhline
\sphinxAtStartPar
\sphinxstylestrong{Test Top\sphinxhyphen{}5}
&
\sphinxAtStartPar
\sphinxstylestrong{94.70 \%}
\\
\sphinxhline
\sphinxAtStartPar
RSRP MAE
&
\sphinxAtStartPar
5.23 dBm
\\
\sphinxbottomrule
\end{tabulary}
\sphinxtableafterendhook\par
\sphinxattableend\end{savenotes}


\subsection{Training Curves}
\label{\detokenize{architecture:id4}}
\sphinxAtStartPar
\sphinxincludegraphics{{6g_training_curves}.png}


\subsection{Confusion Matrix}
\label{\detokenize{architecture:id5}}
\sphinxAtStartPar
\sphinxincludegraphics{{6g_predictive_confusion_matrix}.png}


\subsection{Comparison Plot}
\label{\detokenize{architecture:comparison-plot}}
\sphinxAtStartPar
\sphinxincludegraphics{{exp4_comparison}.png}


\bigskip\hrule\bigskip



\section{Exp 05 — Strategic DeepSet (Context Injection)}
\label{\detokenize{architecture:exp-05-strategic-deepset-context-injection}}\begin{quote}

\sphinxAtStartPar
\sphinxstylestrong{Notebook:} \sphinxcode{\sphinxupquote{notebooks/modeling/05\_strategic\_deepset.ipynb}}\\
\sphinxstylestrong{Checkpoint:} \sphinxcode{\sphinxupquote{models/models\_deprecated/best\_strategic\_deepset.keras}}\\
\sphinxstylestrong{Loss:} \sphinxcode{\sphinxupquote{1.0 × Focal(γ=2.0) + 0.5 × Huber}}
\end{quote}

\sphinxAtStartPar
This experiment injects global contextual state (UE speed, direction cosines, cell load, and one\sphinxhyphen{}hot handover class) into the DeepSet aggregation via concatenation. The goal is to make the model context\sphinxhyphen{}aware of mobility regime and network state beyond just the RF signal values.

\sphinxAtStartPar
\sphinxincludegraphics{{strategic_deepset_architecture}.png}

\begin{sphinxVerbatim}[commandchars=\\\{\}]
Global Context G = [speed, dir\PYGZus{}cos, dir\PYGZus{}sin, cell\PYGZus{}load, ho\PYGZus{}class\PYGZus{}onehot]  ∈ ℝ\PYGZca{}9
    │
    ├─ [Global LSTM (32)] \(\rightarrow\) global embedding gₜ ∈ ℝ\PYGZca{}32
    │
    ├─ [Per\PYGZhy{}Cell LSTM (64)] \(\rightarrow\) φᵢ ∈ ℝ\PYGZca{}64
    │
    ├─ [Concat: φᵢ ‖ gₜ] \(\rightarrow\) enriched cell embedding
    │
    ├─ [RHO MLP (256 \(\rightarrow\) 128)] \(\rightarrow\) aggregated representation
    │
    ├─ [Classifier] \(\rightarrow\) logits ∈ ℝ\PYGZca{}K
    │
    └─ [RSRP Regression Head]
\end{sphinxVerbatim}


\subsection{Results}
\label{\detokenize{architecture:id6}}

\begin{savenotes}\sphinxattablestart
\sphinxthistablewithglobalstyle
\centering
\begin{tabulary}{\linewidth}[t]{TT}
\sphinxtoprule
\sphinxstyletheadfamily 
\sphinxAtStartPar
Metric
&\sphinxstyletheadfamily 
\sphinxAtStartPar
Value
\\
\sphinxmidrule
\sphinxtableatstartofbodyhook
\sphinxAtStartPar
\sphinxstylestrong{Test Top\sphinxhyphen{}1}
&
\sphinxAtStartPar
\sphinxstylestrong{51.35 \%}
\\
\sphinxhline
\sphinxAtStartPar
\sphinxstylestrong{Test Top\sphinxhyphen{}3}
&
\sphinxAtStartPar
\sphinxstylestrong{86.68 \%}
\\
\sphinxhline
\sphinxAtStartPar
\sphinxstylestrong{Test Top\sphinxhyphen{}5}
&
\sphinxAtStartPar
\sphinxstylestrong{95.59 \%}
\\
\sphinxhline
\sphinxAtStartPar
RSRP MAE
&
\sphinxAtStartPar
4.94 dBm
\\
\sphinxbottomrule
\end{tabulary}
\sphinxtableafterendhook\par
\sphinxattableend\end{savenotes}


\subsection{Training Curves}
\label{\detokenize{architecture:id7}}
\sphinxAtStartPar
\sphinxincludegraphics{{strategic_training_curves}.png}


\bigskip\hrule\bigskip



\section{Exp 06 — SOTA Ray\sphinxhyphen{}Tracing Fine\sphinxhyphen{}Tune}
\label{\detokenize{architecture:exp-06-sota-ray-tracing-fine-tune}}\begin{quote}

\sphinxAtStartPar
\sphinxstylestrong{Notebook:} \sphinxcode{\sphinxupquote{notebooks/modeling/06\_sota\_ray\_tracing\_finetuning.ipynb}}\\
\sphinxstylestrong{Source model:} \sphinxcode{\sphinxupquote{models/best\_mtl\_transformer.keras}}\\
\sphinxstylestrong{Fine\sphinxhyphen{}tuned checkpoint:} \sphinxcode{\sphinxupquote{models/best\_sota\_ray\_finetuned.keras}}
\end{quote}

\sphinxAtStartPar
This experiment adapts the best hexagonal\sphinxhyphen{}data model (MTL Transformer) to the high\sphinxhyphen{}fidelity ray\sphinxhyphen{}tracing domain using a \sphinxstylestrong{progressive unfreezing fine\sphinxhyphen{}tuning strategy}:
\begin{enumerate}
\sphinxsetlistlabels{\arabic}{enumi}{enumii}{}{.}%
\item {} 
\sphinxAtStartPar
\sphinxstylestrong{Phase 1 — Head\sphinxhyphen{}only} (25 epochs, LR=1e\sphinxhyphen{}4): freeze encoder, train classifier head only.

\item {} 
\sphinxAtStartPar
\sphinxstylestrong{Phase 2 — Upper layers} (25 epochs, LR=5e\sphinxhyphen{}5): unfreeze upper transformer blocks.

\item {} 
\sphinxAtStartPar
\sphinxstylestrong{Phase 3 — Full model} (15 epochs, LR=1e\sphinxhyphen{}6): end\sphinxhyphen{}to\sphinxhyphen{}end fine\sphinxhyphen{}tuning.

\end{enumerate}

\begin{sphinxVerbatim}[commandchars=\\\{\}]
Base Model (MTL Transformer, frozen encoder)
    │
    Phase 1: train Head only
    Phase 2: unfreeze upper blocks
    Phase 3: full end\PYGZhy{}to\PYGZhy{}end
    │
    \(\rightarrow\) Fine\PYGZhy{}tuned for ray\PYGZhy{}tracing domain
\end{sphinxVerbatim}


\subsection{Results on Ray\sphinxhyphen{}Tracing Test Split}
\label{\detokenize{architecture:results-on-ray-tracing-test-split}}

\begin{savenotes}\sphinxattablestart
\sphinxthistablewithglobalstyle
\centering
\begin{tabulary}{\linewidth}[t]{TTTTT}
\sphinxtoprule
\sphinxstyletheadfamily 
\sphinxAtStartPar
Stage
&\sphinxstyletheadfamily 
\sphinxAtStartPar
Top\sphinxhyphen{}1
&\sphinxstyletheadfamily 
\sphinxAtStartPar
Top\sphinxhyphen{}3
&\sphinxstyletheadfamily 
\sphinxAtStartPar
Top\sphinxhyphen{}5
&\sphinxstyletheadfamily 
\sphinxAtStartPar
RSRP MAE
\\
\sphinxmidrule
\sphinxtableatstartofbodyhook
\sphinxAtStartPar
Baseline (no FT)
&
\sphinxAtStartPar
21.4 \%
&
\sphinxAtStartPar
63.8 \%
&
\sphinxAtStartPar
98.3 \%
&
\sphinxAtStartPar
8.55 dBm
\\
\sphinxhline
\sphinxAtStartPar
\sphinxstylestrong{Fine\sphinxhyphen{}tuned}
&
\sphinxAtStartPar
\sphinxstylestrong{28.1 \%}
&
\sphinxAtStartPar
\sphinxstylestrong{72.5 \%}
&
\sphinxAtStartPar
\sphinxstylestrong{98.6 \%}
&
\sphinxAtStartPar
\sphinxstylestrong{8.11 dBm}
\\
\sphinxbottomrule
\end{tabulary}
\sphinxtableafterendhook\par
\sphinxattableend\end{savenotes}


\subsection{Training Curves}
\label{\detokenize{architecture:id8}}
\sphinxAtStartPar
\sphinxincludegraphics{{ray_finetune_training_curves}.png}


\subsection{Confusion Matrix}
\label{\detokenize{architecture:id9}}
\sphinxAtStartPar
\sphinxincludegraphics{{ray_finetune_confusion_matrix}.png}


\subsection{Occlusion Example}
\label{\detokenize{architecture:occlusion-example}}
\sphinxAtStartPar
\sphinxincludegraphics{{ray_occlusion_example}.png}


\bigskip\hrule\bigskip



\section{Exp 07 — Ray\sphinxhyphen{}ID CrossFormer ★ SOTA / Production}
\label{\detokenize{architecture:exp-07-ray-id-crossformer-sota-production}}\begin{quote}

\sphinxAtStartPar
\sphinxstylestrong{Notebook:} \sphinxcode{\sphinxupquote{notebooks/modeling/07\_ray\_id\_crossformer.ipynb}}\\
\sphinxstylestrong{Production checkpoint:} \sphinxcode{\sphinxupquote{models/best\_ray\_id\_crossformer.keras}}\\
\sphinxstylestrong{Deployment target:} NR edge server (MEC)
\end{quote}

\sphinxAtStartPar
The Ray\sphinxhyphen{}ID CrossFormer is the \sphinxstylestrong{state\sphinxhyphen{}of\sphinxhyphen{}the\sphinxhyphen{}art production model}. It combines:
\begin{itemize}
\item {} 
\sphinxAtStartPar
\sphinxstylestrong{Temporal RF encoding} via GRU (per\sphinxhyphen{}cell, shared weights),

\item {} 
\sphinxAtStartPar
\sphinxstylestrong{Cell identity embeddings} (learned from integer cell IDs, vocabulary size = 172),

\item {} 
\sphinxAtStartPar
\sphinxstylestrong{Serving\sphinxhyphen{}cell flag} embedding (binary: is this cell the current serving cell?),

\item {} 
\sphinxAtStartPar
\sphinxstylestrong{Masked intra\sphinxhyphen{}cell self\sphinxhyphen{}attention} (attend over the candidate set),

\item {} 
\sphinxAtStartPar
\sphinxstylestrong{Horizon cross\sphinxhyphen{}attention} (cross\sphinxhyphen{}attend over the 5 future prediction horizons).

\end{itemize}

\begin{sphinxVerbatim}[commandchars=\\\{\}]
Input:
  ├─ RF features   (K, T=25, 3)   [RSRP, SINR, Load]
  ├─ Cell IDs      (K,)            integer cell identifiers
  └─ Serving flag  (K,)            binary serving\PYGZhy{}cell indicator

Per\PYGZhy{}Cell Temporal Encoder:
  [GRU (96 units, shared weights)] \(\rightarrow\) hᵢ ∈ ℝ\PYGZca{}96

Cell Identity Embedding:
  [Embedding(vocab=172, ID\PYGZus{}DIM=32)] \(\rightarrow\) eᵢ ∈ ℝ\PYGZca{}32

Serving Embedding:
  [Embedding(2, 8)] \(\rightarrow\) sᵢ ∈ ℝ\PYGZca{}8

Fusion:
  [Dense(D\PYGZus{}MODEL=128)] ← concat(hᵢ, eᵢ, sᵢ)

Intra\PYGZhy{}Cell Self\PYGZhy{}Attention:
  [Multi\PYGZhy{}Head Self\PYGZhy{}Attention × 3 blocks, 4 heads, FF=256]  (masked for padding)

Horizon Cross\PYGZhy{}Attention:
  [Learnable horizon query vectors Q ∈ ℝ\PYGZca{}(H×D)]
  [Cross\PYGZhy{}Attention over cell keys/values]
  \(\rightarrow\) per\PYGZhy{}horizon logits ∈ ℝ\PYGZca{}(H×K)

Output:
  Softmax probabilities per horizon per cell  \(\rightarrow\)  shape (H=5, K)
\end{sphinxVerbatim}


\subsection{Hyperparameters}
\label{\detokenize{architecture:id10}}

\begin{savenotes}\sphinxattablestart
\sphinxthistablewithglobalstyle
\centering
\begin{tabulary}{\linewidth}[t]{TT}
\sphinxtoprule
\sphinxstyletheadfamily 
\sphinxAtStartPar
Parameter
&\sphinxstyletheadfamily 
\sphinxAtStartPar
Value
\\
\sphinxmidrule
\sphinxtableatstartofbodyhook
\sphinxAtStartPar
\sphinxcode{\sphinxupquote{D\_MODEL}}
&
\sphinxAtStartPar
128
\\
\sphinxhline
\sphinxAtStartPar
\sphinxcode{\sphinxupquote{ID\_DIM}}
&
\sphinxAtStartPar
32
\\
\sphinxhline
\sphinxAtStartPar
\sphinxcode{\sphinxupquote{GRU\_UNITS}}
&
\sphinxAtStartPar
96
\\
\sphinxhline
\sphinxAtStartPar
\sphinxcode{\sphinxupquote{N\_HEADS}}
&
\sphinxAtStartPar
4
\\
\sphinxhline
\sphinxAtStartPar
\sphinxcode{\sphinxupquote{FF\_DIM}}
&
\sphinxAtStartPar
256
\\
\sphinxhline
\sphinxAtStartPar
\sphinxcode{\sphinxupquote{N\_BLOCKS}}
&
\sphinxAtStartPar
3
\\
\sphinxhline
\sphinxAtStartPar
\sphinxcode{\sphinxupquote{DROPOUT}}
&
\sphinxAtStartPar
0.18
\\
\sphinxhline
\sphinxAtStartPar
\sphinxcode{\sphinxupquote{FOCAL\_GAMMA}}
&
\sphinxAtStartPar
1.5
\\
\sphinxhline
\sphinxAtStartPar
\sphinxcode{\sphinxupquote{FOCAL\_ALPHA}}
&
\sphinxAtStartPar
0.75
\\
\sphinxhline
\sphinxAtStartPar
\sphinxcode{\sphinxupquote{LABEL\_SMOOTHING}}
&
\sphinxAtStartPar
0.01
\\
\sphinxhline
\sphinxAtStartPar
\sphinxcode{\sphinxupquote{BATCH\_SIZE}}
&
\sphinxAtStartPar
64
\\
\sphinxhline
\sphinxAtStartPar
\sphinxcode{\sphinxupquote{EPOCHS}}
&
\sphinxAtStartPar
90
\\
\sphinxhline
\sphinxAtStartPar
\sphinxcode{\sphinxupquote{LR}}
&
\sphinxAtStartPar
3e\sphinxhyphen{}4
\\
\sphinxhline
\sphinxAtStartPar
\sphinxcode{\sphinxupquote{MIN\_LR}}
&
\sphinxAtStartPar
2e\sphinxhyphen{}6
\\
\sphinxhline
\sphinxAtStartPar
\sphinxcode{\sphinxupquote{PATIENCE}}
&
\sphinxAtStartPar
14
\\
\sphinxbottomrule
\end{tabulary}
\sphinxtableafterendhook\par
\sphinxattableend\end{savenotes}


\subsection{Results on Ray\sphinxhyphen{}Tracing Domain}
\label{\detokenize{architecture:results-on-ray-tracing-domain}}

\begin{savenotes}\sphinxattablestart
\sphinxthistablewithglobalstyle
\centering
\begin{tabulary}{\linewidth}[t]{TTTT}
\sphinxtoprule
\sphinxstyletheadfamily 
\sphinxAtStartPar
Split
&\sphinxstyletheadfamily 
\sphinxAtStartPar
Top\sphinxhyphen{}1
&\sphinxstyletheadfamily 
\sphinxAtStartPar
Top\sphinxhyphen{}3
&\sphinxstyletheadfamily 
\sphinxAtStartPar
Top\sphinxhyphen{}5
\\
\sphinxmidrule
\sphinxtableatstartofbodyhook
\sphinxAtStartPar
\sphinxstylestrong{Validation}
&
\sphinxAtStartPar
\sphinxstylestrong{59.96 \%}
&
\sphinxAtStartPar
\sphinxstylestrong{90.86 \%}
&
\sphinxAtStartPar
\sphinxstylestrong{98.80 \%}
\\
\sphinxhline
\sphinxAtStartPar
\sphinxstylestrong{Test}
&
\sphinxAtStartPar
\sphinxstylestrong{57.60 \%}
&
\sphinxAtStartPar
\sphinxstylestrong{90.87 \%}
&
\sphinxAtStartPar
\sphinxstylestrong{98.87 \%}
\\
\sphinxbottomrule
\end{tabulary}
\sphinxtableafterendhook\par
\sphinxattableend\end{savenotes}
\begin{quote}

\sphinxAtStartPar
\sphinxstylestrong{Top\sphinxhyphen{}3 accuracy (90.87 \%) is the primary production KPI} — in a live NR network, the RAN can negotiate among the top\sphinxhyphen{}3 predicted candidates using X2/Xn signaling, maximising flexibility while dramatically reducing RLF risk. This is the standard operational metric in edge\sphinxhyphen{}AI handover pipelines.
\end{quote}


\subsection{Training Curves}
\label{\detokenize{architecture:id11}}
\sphinxAtStartPar
\sphinxincludegraphics{{ray_id_crossformer_training_curves}.png}


\subsection{Confusion Matrix}
\label{\detokenize{architecture:id12}}
\sphinxAtStartPar
\sphinxincludegraphics{{ray_id_crossformer_confusion_matrix}.png}


\bigskip\hrule\bigskip



\section{Model Comparison Summary}
\label{\detokenize{architecture:model-comparison-summary}}

\begin{savenotes}\sphinxattablestart
\sphinxthistablewithglobalstyle
\centering
\begin{tabulary}{\linewidth}[t]{TTTTTT}
\sphinxtoprule
\sphinxstyletheadfamily 
\sphinxAtStartPar
Model
&\sphinxstyletheadfamily 
\sphinxAtStartPar
Task
&\sphinxstyletheadfamily 
\sphinxAtStartPar
Top\sphinxhyphen{}1
&\sphinxstyletheadfamily 
\sphinxAtStartPar
Top\sphinxhyphen{}3
&\sphinxstyletheadfamily 
\sphinxAtStartPar
Top\sphinxhyphen{}5
&\sphinxstyletheadfamily 
\sphinxAtStartPar
RSRP MAE
\\
\sphinxmidrule
\sphinxtableatstartofbodyhook
\sphinxAtStartPar
Temporal DeepSet (Exp 01)
&
\sphinxAtStartPar
\sphinxstylestrong{Single\sphinxhyphen{}step (baseline only)}
&
\sphinxAtStartPar
99.82 \%
&
\sphinxAtStartPar
100.0 \%
&
\sphinxAtStartPar
100.0 \%
&
\sphinxAtStartPar
—
\\
\sphinxhline
\sphinxAtStartPar
SpatioTemporal DeepSet (Exp 02)
&
\sphinxAtStartPar
5\sphinxhyphen{}step
&
\sphinxAtStartPar
56.13 \% (avg)
&
\sphinxAtStartPar
—
&
\sphinxAtStartPar
—
&
\sphinxAtStartPar
5.56 dBm
\\
\sphinxhline
\sphinxAtStartPar
MTL Transformer (Exp 03)
&
\sphinxAtStartPar
5\sphinxhyphen{}step
&
\sphinxAtStartPar
55.79 \%
&
\sphinxAtStartPar
86.02 \%
&
\sphinxAtStartPar
94.99 \%
&
\sphinxAtStartPar
4.20 dBm
\\
\sphinxhline
\sphinxAtStartPar
6G Predictive (Exp 04)
&
\sphinxAtStartPar
5\sphinxhyphen{}step
&
\sphinxAtStartPar
54.93 \%
&
\sphinxAtStartPar
85.47 \%
&
\sphinxAtStartPar
94.70 \%
&
\sphinxAtStartPar
5.23 dBm
\\
\sphinxhline
\sphinxAtStartPar
Strategic DeepSet (Exp 05)
&
\sphinxAtStartPar
5\sphinxhyphen{}step
&
\sphinxAtStartPar
51.35 \%
&
\sphinxAtStartPar
86.68 \%
&
\sphinxAtStartPar
95.59 \%
&
\sphinxAtStartPar
4.94 dBm
\\
\sphinxhline
\sphinxAtStartPar
Ray FT (Exp 06)
&
\sphinxAtStartPar
5\sphinxhyphen{}step (ray domain)
&
\sphinxAtStartPar
28.12 \%
&
\sphinxAtStartPar
72.55 \%
&
\sphinxAtStartPar
98.57 \%
&
\sphinxAtStartPar
8.11 dBm
\\
\sphinxhline
\sphinxAtStartPar
\sphinxstylestrong{Ray\sphinxhyphen{}ID CrossFormer (Exp 07) ★}
&
\sphinxAtStartPar
\sphinxstylestrong{5\sphinxhyphen{}step (ray domain)}
&
\sphinxAtStartPar
\sphinxstylestrong{57.60 \%}
&
\sphinxAtStartPar
\sphinxstylestrong{90.87 \%}
&
\sphinxAtStartPar
\sphinxstylestrong{98.87 \%}
&
\sphinxAtStartPar
—
\\
\sphinxbottomrule
\end{tabulary}
\sphinxtableafterendhook\par
\sphinxattableend\end{savenotes}
\begin{quote}

\sphinxAtStartPar
\sphinxstylestrong{Note on Top\sphinxhyphen{}3 as Production KPI:} In real NR deployments, the network selects from the top\sphinxhyphen{}3 candidate cells based on live X2/Xn signaling, QoS constraints, and load balancing. A Top\sphinxhyphen{}3 accuracy of \sphinxstylestrong{90.87 \%} means that in 9 out of 10 handover events, the optimal target cell is within the model’s top\sphinxhyphen{}3 recommendations — ensuring maximum operational flexibility and near\sphinxhyphen{}elimination of blind\sphinxhyphen{}handover failures.
\end{quote}

\sphinxstepscope


\chapter{Training Reference}
\label{\detokenize{training:training-reference}}\label{\detokenize{training::doc}}
\sphinxAtStartPar
This page describes the full training methodology applied across the experimentation pipeline: data preprocessing, loss functions, learning rate scheduling, evaluation protocol, and per\sphinxhyphen{}experiment training specifics. All \sphinxstylestrong{production experiments (Exp 02\textendash{}07) target 5\sphinxhyphen{}step future horizon prediction}. Only Exp 01 is a single\sphinxhyphen{}step baseline.


\bigskip\hrule\bigskip



\section{1. Data Pipeline}
\label{\detokenize{training:data-pipeline}}

\subsection{1.1 Raw Data Sources}
\label{\detokenize{training:raw-data-sources}}

\begin{savenotes}\sphinxattablestart
\sphinxthistablewithglobalstyle
\centering
\begin{tabulary}{\linewidth}[t]{TTT}
\sphinxtoprule
\sphinxstyletheadfamily 
\sphinxAtStartPar
Dataset
&\sphinxstyletheadfamily 
\sphinxAtStartPar
Path
&\sphinxstyletheadfamily 
\sphinxAtStartPar
Description
\\
\sphinxmidrule
\sphinxtableatstartofbodyhook
\sphinxAtStartPar
Hexagonal simulation
&
\sphinxAtStartPar
\sphinxcode{\sphinxupquote{dataset/raw/handover\_dataset.csv}}
&
\sphinxAtStartPar
Large\sphinxhyphen{}scale UMi channel model dataset (primary training corpus)
\\
\sphinxhline
\sphinxAtStartPar
Ray\sphinxhyphen{}tracing simulation
&
\sphinxAtStartPar
\sphinxcode{\sphinxupquote{dataset/raw/handover\_dataset\_ray\_tracing.csv}}
&
\sphinxAtStartPar
High\sphinxhyphen{}fidelity ray\sphinxhyphen{}tracing dataset (fine\sphinxhyphen{}tuning / Exp 06\textendash{}07)
\\
\sphinxbottomrule
\end{tabulary}
\sphinxtableafterendhook\par
\sphinxattableend\end{savenotes}


\subsection{1.2 UE\sphinxhyphen{}Level Train/Val/Test Split}
\label{\detokenize{training:ue-level-train-val-test-split}}
\sphinxAtStartPar
All experiments use a \sphinxstylestrong{UE\sphinxhyphen{}level split} loaded from \sphinxcode{\sphinxupquote{dataset/processed/ue\_split.json}}. This prevents temporal leakage — a UE’s data appears in exactly one partition. Typical split ratios are 70/15/15.

\begin{sphinxVerbatim}[commandchars=\\\{\}]
\PYG{c+c1}{\PYGZsh{} UE\PYGZhy{}level split prevents time leakage across train/val/test}
\PYG{k}{with} \PYG{n+nb}{open}\PYG{p}{(}\PYG{l+s+s2}{\PYGZdq{}}\PYG{l+s+s2}{dataset/processed/ue\PYGZus{}split.json}\PYG{l+s+s2}{\PYGZdq{}}\PYG{p}{)} \PYG{k}{as} \PYG{n}{f}\PYG{p}{:}
    \PYG{n}{ue\PYGZus{}split} \PYG{o}{=} \PYG{n}{json}\PYG{o}{.}\PYG{n}{load}\PYG{p}{(}\PYG{n}{f}\PYG{p}{)}
    \PYG{n}{train\PYGZus{}ues} \PYG{o}{=} \PYG{n}{ue\PYGZus{}split}\PYG{p}{[}\PYG{l+s+s2}{\PYGZdq{}}\PYG{l+s+s2}{train}\PYG{l+s+s2}{\PYGZdq{}}\PYG{p}{]}   \PYG{c+c1}{\PYGZsh{} \PYGZti{}70\PYGZpc{} of UE trajectories}
    \PYG{n}{val\PYGZus{}ues}   \PYG{o}{=} \PYG{n}{ue\PYGZus{}split}\PYG{p}{[}\PYG{l+s+s2}{\PYGZdq{}}\PYG{l+s+s2}{val}\PYG{l+s+s2}{\PYGZdq{}}\PYG{p}{]}     \PYG{c+c1}{\PYGZsh{} \PYGZti{}15\PYGZpc{}}
    \PYG{n}{test\PYGZus{}ues}  \PYG{o}{=} \PYG{n}{ue\PYGZus{}split}\PYG{p}{[}\PYG{l+s+s2}{\PYGZdq{}}\PYG{l+s+s2}{test}\PYG{l+s+s2}{\PYGZdq{}}\PYG{p}{]}    \PYG{c+c1}{\PYGZsh{} \PYGZti{}15\PYGZpc{}  ← never seen during training}
\end{sphinxVerbatim}


\subsection{1.3 Sliding Window Construction}
\label{\detokenize{training:sliding-window-construction}}
\sphinxAtStartPar
Each sample is a sliding window of \sphinxcode{\sphinxupquote{OBS\_STEPS}} consecutive timesteps (sampled every 200 ms):

\begin{sphinxVerbatim}[commandchars=\\\{\}]
[ t\PYGZhy{}(T\PYGZhy{}1)  ...  t\PYGZhy{}1  t ]  ──►  Predict optimal cell at  [ t+1, t+2, t+3, t+4, t+5 ]
\end{sphinxVerbatim}

\sphinxAtStartPar
The window shape is \sphinxcode{\sphinxupquote{(K, T, F)}}:
\begin{itemize}
\item {} 
\sphinxAtStartPar
\sphinxcode{\sphinxupquote{K = MAX\_CELLS = 10}} — padded candidate cell set

\item {} 
\sphinxAtStartPar
\sphinxcode{\sphinxupquote{T = OBS\_STEPS}} — 25 or 200 timesteps depending on experiment

\item {} 
\sphinxAtStartPar
\sphinxcode{\sphinxupquote{F}} — feature count (3\textendash{}5 depending on experiment)

\end{itemize}


\subsection{1.4 Neighbor\sphinxhyphen{}Axis Shuffling (Anti\sphinxhyphen{}Leakage)}
\label{\detokenize{training:neighbor-axis-shuffling-anti-leakage}}
\sphinxAtStartPar
A critical anti\sphinxhyphen{}leakage measure is applied to \sphinxstylestrong{every sample}: the \sphinxcode{\sphinxupquote{K}} candidate cells are randomly permuted before caching, and the label is remapped accordingly. This prevents the model from learning the trivial shortcut of “index 0 is always optimal” (which was an artefact of score\sphinxhyphen{}sorted raw CSVs).

\begin{sphinxVerbatim}[commandchars=\\\{\}]
\PYG{c+c1}{\PYGZsh{} Applied per window in all production experiments}
\PYG{n}{perm}  \PYG{o}{=} \PYG{n}{rng}\PYG{o}{.}\PYG{n}{permutation}\PYG{p}{(}\PYG{n}{MAX\PYGZus{}CELLS}\PYG{p}{)}
\PYG{n}{X}     \PYG{o}{=} \PYG{n}{X}\PYG{p}{[}\PYG{n}{perm}\PYG{p}{]}          \PYG{c+c1}{\PYGZsh{} (K, T, F) shuffled}
\PYG{n}{label} \PYG{o}{=} \PYG{n}{inv\PYGZus{}perm}\PYG{p}{[}\PYG{n}{label}\PYG{p}{]}  \PYG{c+c1}{\PYGZsh{} label index remapped after shuffle}
\end{sphinxVerbatim}


\subsection{1.5 Feature Schema}
\label{\detokenize{training:feature-schema}}

\begin{savenotes}\sphinxattablestart
\sphinxthistablewithglobalstyle
\centering
\begin{tabulary}{\linewidth}[t]{TTT}
\sphinxtoprule
\sphinxstyletheadfamily 
\sphinxAtStartPar
Feature
&\sphinxstyletheadfamily 
\sphinxAtStartPar
Description
&\sphinxstyletheadfamily 
\sphinxAtStartPar
Included in
\\
\sphinxmidrule
\sphinxtableatstartofbodyhook
\sphinxAtStartPar
\sphinxcode{\sphinxupquote{nb\_rsrp}}
&
\sphinxAtStartPar
Neighbor cell RSRP (dBm)
&
\sphinxAtStartPar
All experiments
\\
\sphinxhline
\sphinxAtStartPar
\sphinxcode{\sphinxupquote{nb\_sinr}}
&
\sphinxAtStartPar
Neighbor cell SINR (dB)
&
\sphinxAtStartPar
All experiments
\\
\sphinxhline
\sphinxAtStartPar
\sphinxcode{\sphinxupquote{nb\_load}}
&
\sphinxAtStartPar
Neighbor cell load (0\textendash{}1)
&
\sphinxAtStartPar
All experiments
\\
\sphinxhline
\sphinxAtStartPar
\sphinxcode{\sphinxupquote{zero\_anchor}}
&
\sphinxAtStartPar
LTE anchor serving\sphinxhyphen{}flag
&
\sphinxAtStartPar
Exp 03, 04, 06
\\
\sphinxhline
\sphinxAtStartPar
\sphinxcode{\sphinxupquote{speed / dir\_cos / dir\_sin}}
&
\sphinxAtStartPar
UE mobility state
&
\sphinxAtStartPar
Exp 05
\\
\sphinxhline
\sphinxAtStartPar
Cell ID integers
&
\sphinxAtStartPar
Discrete cell identifiers
&
\sphinxAtStartPar
Exp 07 (ID embedding)
\\
\sphinxbottomrule
\end{tabulary}
\sphinxtableafterendhook\par
\sphinxattableend\end{savenotes}


\bigskip\hrule\bigskip



\section{2. Loss Functions}
\label{\detokenize{training:loss-functions}}

\subsection{2.1 Focal Loss}
\label{\detokenize{training:focal-loss}}
\sphinxAtStartPar
All classification heads use Focal Loss to handle class imbalance across the 172 candidate cells:
\begin{equation*}
\begin{split}\mathcal{L}_{\text{focal}}(p_t) = -\alpha_t (1 - p_t)^\gamma \log(p_t)\end{split}
\end{equation*}
\sphinxAtStartPar
Where:
\begin{itemize}
\item {} 
\sphinxAtStartPar
\(\alpha\) focuses on the minority class (positive handover events)

\item {} 
\sphinxAtStartPar
\(\gamma\) down\sphinxhyphen{}weights easy (well\sphinxhyphen{}classified) samples

\end{itemize}


\begin{savenotes}\sphinxattablestart
\sphinxthistablewithglobalstyle
\centering
\begin{tabulary}{\linewidth}[t]{TTT}
\sphinxtoprule
\sphinxstyletheadfamily 
\sphinxAtStartPar
Experiment
&\sphinxstyletheadfamily 
\sphinxAtStartPar
γ
&\sphinxstyletheadfamily 
\sphinxAtStartPar
α
\\
\sphinxmidrule
\sphinxtableatstartofbodyhook
\sphinxAtStartPar
Exp 01\textendash{}05
&
\sphinxAtStartPar
2.0
&
\sphinxAtStartPar
0.25
\\
\sphinxhline
\sphinxAtStartPar
Exp 06 (fine\sphinxhyphen{}tune)
&
\sphinxAtStartPar
2.0
&
\sphinxAtStartPar
0.25
\\
\sphinxhline
\sphinxAtStartPar
\sphinxstylestrong{Exp 07 (CrossFormer)}
&
\sphinxAtStartPar
\sphinxstylestrong{1.5}
&
\sphinxAtStartPar
\sphinxstylestrong{0.75}
\\
\sphinxbottomrule
\end{tabulary}
\sphinxtableafterendhook\par
\sphinxattableend\end{savenotes}


\subsection{2.2 Auxiliary Regression Loss (Huber / MSE)}
\label{\detokenize{training:auxiliary-regression-loss-huber-mse}}
\sphinxAtStartPar
Multi\sphinxhyphen{}task experiments (03, 04, 05) add a signal regression head with a Huber or MSE loss:
\begin{equation*}
\begin{split}\mathcal{L}_{\text{total}} = \lambda_{\text{cls}} \cdot \mathcal{L}_{\text{focal}} + \lambda_{\text{reg}} \cdot \mathcal{L}_{\text{huber}}\end{split}
\end{equation*}

\begin{savenotes}\sphinxattablestart
\sphinxthistablewithglobalstyle
\centering
\begin{tabulary}{\linewidth}[t]{TTT}
\sphinxtoprule
\sphinxstyletheadfamily 
\sphinxAtStartPar
Experiment
&\sphinxstyletheadfamily 
\sphinxAtStartPar
λ\_cls
&\sphinxstyletheadfamily 
\sphinxAtStartPar
λ\_reg
\\
\sphinxmidrule
\sphinxtableatstartofbodyhook
\sphinxAtStartPar
Exp 03
&
\sphinxAtStartPar
1.0
&
\sphinxAtStartPar
0.5
\\
\sphinxhline
\sphinxAtStartPar
Exp 04
&
\sphinxAtStartPar
1.0
&
\sphinxAtStartPar
0.7
\\
\sphinxhline
\sphinxAtStartPar
Exp 05
&
\sphinxAtStartPar
1.0
&
\sphinxAtStartPar
0.5
\\
\sphinxbottomrule
\end{tabulary}
\sphinxtableafterendhook\par
\sphinxattableend\end{savenotes}


\subsection{2.3 Label Smoothing}
\label{\detokenize{training:label-smoothing}}
\sphinxAtStartPar
Applied to the classification target to regularise overconfident predictions:


\begin{savenotes}\sphinxattablestart
\sphinxthistablewithglobalstyle
\centering
\begin{tabulary}{\linewidth}[t]{TT}
\sphinxtoprule
\sphinxstyletheadfamily 
\sphinxAtStartPar
Experiment
&\sphinxstyletheadfamily 
\sphinxAtStartPar
Label Smoothing
\\
\sphinxmidrule
\sphinxtableatstartofbodyhook
\sphinxAtStartPar
Exp 01
&
\sphinxAtStartPar
0.10
\\
\sphinxhline
\sphinxAtStartPar
Exp 02
&
\sphinxAtStartPar
0.10
\\
\sphinxhline
\sphinxAtStartPar
Exp 06 (fine\sphinxhyphen{}tune)
&
\sphinxAtStartPar
0.02
\\
\sphinxhline
\sphinxAtStartPar
\sphinxstylestrong{Exp 07 (CrossFormer)}
&
\sphinxAtStartPar
\sphinxstylestrong{0.01}
\\
\sphinxbottomrule
\end{tabulary}
\sphinxtableafterendhook\par
\sphinxattableend\end{savenotes}


\bigskip\hrule\bigskip



\section{3. Learning Rate Scheduling}
\label{\detokenize{training:learning-rate-scheduling}}

\subsection{3.1 Warmup + Cosine Decay (Exp 01\textendash{}06)}
\label{\detokenize{training:warmup-cosine-decay-exp-0106}}
\sphinxAtStartPar
A linear warmup phase is followed by cosine decay, providing stable early\sphinxhyphen{}epoch convergence and smooth annealing:

\begin{sphinxVerbatim}[commandchars=\\\{\}]
Epochs 0 \(\rightarrow\) LR\PYGZus{}WARMUP\PYGZus{}EP  : linear ramp from 0 \(\rightarrow\) LR\PYGZus{}INIT
Epochs LR\PYGZus{}WARMUP\PYGZus{}EP \(\rightarrow\) LR\PYGZus{}DECAY\PYGZus{}EP : cosine decay to LR\PYGZus{}MIN
Epochs LR\PYGZus{}DECAY\PYGZus{}EP \(\rightarrow\) EPOCHS : flat at LR\PYGZus{}MIN
\end{sphinxVerbatim}


\begin{savenotes}\sphinxattablestart
\sphinxthistablewithglobalstyle
\centering
\begin{tabulary}{\linewidth}[t]{TTTTT}
\sphinxtoprule
\sphinxstyletheadfamily 
\sphinxAtStartPar
Experiment
&\sphinxstyletheadfamily 
\sphinxAtStartPar
LR\_INIT
&\sphinxstyletheadfamily 
\sphinxAtStartPar
Warmup Epochs
&\sphinxstyletheadfamily 
\sphinxAtStartPar
Decay Epochs
&\sphinxstyletheadfamily 
\sphinxAtStartPar
Total Epochs
\\
\sphinxmidrule
\sphinxtableatstartofbodyhook
\sphinxAtStartPar
Exp 01
&
\sphinxAtStartPar
1e\sphinxhyphen{}3
&
\sphinxAtStartPar
4
&
\sphinxAtStartPar
20
&
\sphinxAtStartPar
60
\\
\sphinxhline
\sphinxAtStartPar
Exp 02
&
\sphinxAtStartPar
1e\sphinxhyphen{}3
&
\sphinxAtStartPar
4
&
\sphinxAtStartPar
20
&
\sphinxAtStartPar
60
\\
\sphinxhline
\sphinxAtStartPar
Exp 03
&
\sphinxAtStartPar
1e\sphinxhyphen{}3
&
\sphinxAtStartPar
4
&
\sphinxAtStartPar
20
&
\sphinxAtStartPar
60
\\
\sphinxhline
\sphinxAtStartPar
Exp 04 (Optuna)
&
\sphinxAtStartPar
1.68e\sphinxhyphen{}3
&
\sphinxAtStartPar
5
&
\sphinxAtStartPar
25
&
\sphinxAtStartPar
60
\\
\sphinxhline
\sphinxAtStartPar
Exp 05
&
\sphinxAtStartPar
5e\sphinxhyphen{}4
&
\sphinxAtStartPar
4
&
\sphinxAtStartPar
20
&
\sphinxAtStartPar
60
\\
\sphinxbottomrule
\end{tabulary}
\sphinxtableafterendhook\par
\sphinxattableend\end{savenotes}


\subsection{3.2 ReduceLROnPlateau + EarlyStopping (Exp 07 — CrossFormer)}
\label{\detokenize{training:reducelronplateau-earlystopping-exp-07-crossformer}}
\sphinxAtStartPar
The Ray\sphinxhyphen{}ID CrossFormer uses a simpler ReduceLROnPlateau callback with early stopping, more appropriate for the smaller ray\sphinxhyphen{}tracing dataset:

\begin{sphinxVerbatim}[commandchars=\\\{\}]
\PYG{n}{ReduceLROnPlateau}\PYG{p}{(}\PYG{n}{monitor}\PYG{o}{=}\PYG{l+s+s1}{\PYGZsq{}}\PYG{l+s+s1}{val\PYGZus{}loss}\PYG{l+s+s1}{\PYGZsq{}}\PYG{p}{,} \PYG{n}{factor}\PYG{o}{=}\PYG{l+m+mf}{0.5}\PYG{p}{,} \PYG{n}{patience}\PYG{o}{=}\PYG{l+m+mi}{7}\PYG{p}{,} \PYG{n}{min\PYGZus{}lr}\PYG{o}{=}\PYG{l+m+mf}{2e\PYGZhy{}6}\PYG{p}{)}
\PYG{n}{EarlyStopping}\PYG{p}{(}\PYG{n}{monitor}\PYG{o}{=}\PYG{l+s+s1}{\PYGZsq{}}\PYG{l+s+s1}{val\PYGZus{}top3\PYGZus{}accuracy}\PYG{l+s+s1}{\PYGZsq{}}\PYG{p}{,} \PYG{n}{patience}\PYG{o}{=}\PYG{l+m+mi}{14}\PYG{p}{,} \PYG{n}{restore\PYGZus{}best\PYGZus{}weights}\PYG{o}{=}\PYG{k+kc}{True}\PYG{p}{)}
\end{sphinxVerbatim}


\begin{savenotes}\sphinxattablestart
\sphinxthistablewithglobalstyle
\centering
\begin{tabulary}{\linewidth}[t]{TT}
\sphinxtoprule
\sphinxstyletheadfamily 
\sphinxAtStartPar
Parameter
&\sphinxstyletheadfamily 
\sphinxAtStartPar
Value
\\
\sphinxmidrule
\sphinxtableatstartofbodyhook
\sphinxAtStartPar
Initial LR
&
\sphinxAtStartPar
3e\sphinxhyphen{}4
\\
\sphinxhline
\sphinxAtStartPar
Min LR
&
\sphinxAtStartPar
2e\sphinxhyphen{}6
\\
\sphinxhline
\sphinxAtStartPar
Patience (early stop)
&
\sphinxAtStartPar
14
\\
\sphinxhline
\sphinxAtStartPar
Total epochs
&
\sphinxAtStartPar
90
\\
\sphinxbottomrule
\end{tabulary}
\sphinxtableafterendhook\par
\sphinxattableend\end{savenotes}


\bigskip\hrule\bigskip



\section{4. Training Infrastructure}
\label{\detokenize{training:training-infrastructure}}

\subsection{4.1 Hardware \& Framework}
\label{\detokenize{training:hardware-framework}}

\begin{savenotes}\sphinxattablestart
\sphinxthistablewithglobalstyle
\centering
\begin{tabulary}{\linewidth}[t]{TT}
\sphinxtoprule
\sphinxstyletheadfamily 
\sphinxAtStartPar
Item
&\sphinxstyletheadfamily 
\sphinxAtStartPar
Detail
\\
\sphinxmidrule
\sphinxtableatstartofbodyhook
\sphinxAtStartPar
Framework
&
\sphinxAtStartPar
Keras 3 / TensorFlow 2.x backend
\\
\sphinxhline
\sphinxAtStartPar
Mixed precision
&
\sphinxAtStartPar
\sphinxcode{\sphinxupquote{float16}} compute, \sphinxcode{\sphinxupquote{float32}} weights (GPU)
\\
\sphinxhline
\sphinxAtStartPar
Dataset API
&
\sphinxAtStartPar
\sphinxcode{\sphinxupquote{tf.data.Dataset}} pipeline pinned to CPU
\\
\sphinxhline
\sphinxAtStartPar
Experiment tracking
&
\sphinxAtStartPar
MLflow (\sphinxcode{\sphinxupquote{mlflow/mlruns/}})
\\
\sphinxhline
\sphinxAtStartPar
TensorBoard logs
&
\sphinxAtStartPar
\sphinxcode{\sphinxupquote{tb\_logs/\textless{}experiment\textgreater{}/}}
\\
\sphinxbottomrule
\end{tabulary}
\sphinxtableafterendhook\par
\sphinxattableend\end{savenotes}


\subsection{4.2 Dataset Caching}
\label{\detokenize{training:dataset-caching}}
\sphinxAtStartPar
All experiments pre\sphinxhyphen{}build a numpy cache to avoid repeated windowing:

\begin{sphinxVerbatim}[commandchars=\\\{\}]
dataset/
├── temporal\PYGZus{}deepset\PYGZus{}cache/          \PYGZsh{} Exp 01
├── temporal\PYGZus{}deepset\PYGZus{}production\PYGZus{}cache/ \PYGZsh{} Exp 02
├── mtl\PYGZus{}cache/                       \PYGZsh{} Exp 03
├── 6g\PYGZus{}cache/                        \PYGZsh{} Exp 04
├── strategic\PYGZus{}cache/                 \PYGZsh{} Exp 05
├── ray\PYGZus{}tracing\PYGZus{}finetune\PYGZus{}cache/      \PYGZsh{} Exp 06\textendash{}07
└── processed/ue\PYGZus{}split.json          \PYGZsh{} shared UE split
\end{sphinxVerbatim}


\subsection{4.3 Callbacks}
\label{\detokenize{training:callbacks}}
\sphinxAtStartPar
Every experiment uses the following standard callback set:

\begin{sphinxVerbatim}[commandchars=\\\{\}]
\PYG{n}{ModelCheckpoint}\PYG{p}{(}\PYG{n}{filepath}\PYG{o}{=}\PYG{l+s+s1}{\PYGZsq{}}\PYG{l+s+s1}{models/best\PYGZus{}\PYGZlt{}exp\PYGZgt{}.keras}\PYG{l+s+s1}{\PYGZsq{}}\PYG{p}{,}
                \PYG{n}{monitor}\PYG{o}{=}\PYG{l+s+s1}{\PYGZsq{}}\PYG{l+s+s1}{val\PYGZus{}top3\PYGZus{}accuracy}\PYG{l+s+s1}{\PYGZsq{}}\PYG{p}{,}    \PYG{c+c1}{\PYGZsh{} primary monitor}
                \PYG{n}{save\PYGZus{}best\PYGZus{}only}\PYG{o}{=}\PYG{k+kc}{True}\PYG{p}{)}
\PYG{n}{TensorBoard}\PYG{p}{(}\PYG{n}{log\PYGZus{}dir}\PYG{o}{=}\PYG{l+s+s1}{\PYGZsq{}}\PYG{l+s+s1}{tb\PYGZus{}logs/\PYGZlt{}exp\PYGZgt{}/}\PYG{l+s+s1}{\PYGZsq{}}\PYG{p}{)}
\PYG{n}{CSVLogger}\PYG{p}{(}\PYG{l+s+s1}{\PYGZsq{}}\PYG{l+s+s1}{metrics/\PYGZlt{}exp\PYGZgt{}/training\PYGZus{}log.csv}\PYG{l+s+s1}{\PYGZsq{}}\PYG{p}{)}
\end{sphinxVerbatim}
\begin{quote}

\sphinxAtStartPar
\sphinxstylestrong{\sphinxcode{\sphinxupquote{val\_top3\_accuracy}} as the checkpoint monitor} — Top\sphinxhyphen{}3 accuracy is used as the primary validation monitor because it is the most operationally relevant metric in production. The NR network uses the top\sphinxhyphen{}3 predicted candidates for X2/Xn preparation, so optimising for Top\sphinxhyphen{}3 directly aligns training with deployment objectives.
\end{quote}


\bigskip\hrule\bigskip



\section{5. Experiment Training Details}
\label{\detokenize{training:experiment-training-details}}

\subsection{Exp 01 — Temporal DeepSet (Single\sphinxhyphen{}Step Baseline)}
\label{\detokenize{training:exp-01-temporal-deepset-single-step-baseline}}\begin{quote}

\sphinxAtStartPar
⚠️ This is the \sphinxstylestrong{single\sphinxhyphen{}step selection} experiment. The model predicts the best current cell. It is used as a \sphinxstylestrong{proof\sphinxhyphen{}of\sphinxhyphen{}concept baseline} and was the first modeling experimentation in this project. All subsequent experiments solve the harder 5\sphinxhyphen{}step future prediction problem.
\end{quote}

\begin{sphinxVerbatim}[commandchars=\\\{\}]
conda\PYG{+w}{ }run\PYG{+w}{ }\PYGZhy{}n\PYG{+w}{ }tda\PYGZhy{}handover\PYG{+w}{ }python\PYG{+w}{ }src/train\PYGZus{}model.py\PYG{+w}{ }\PYGZhy{}\PYGZhy{}label\PYG{+w}{ }optimal
\end{sphinxVerbatim}

\sphinxAtStartPar
\sphinxincludegraphics{{temporal_deepset_training_curves}.png}


\subsection{Exp 02 — SpatioTemporal DeepSet Production}
\label{\detokenize{training:exp-02-spatiotemporal-deepset-production}}
\sphinxAtStartPar
First transition to \sphinxstylestrong{5\sphinxhyphen{}step multi\sphinxhyphen{}horizon} prediction. Uses a long 200\sphinxhyphen{}timestep observation window (40 seconds of history at 200 ms resolution) and a multi\sphinxhyphen{}task loss combining future cell classification and signal regression.

\sphinxAtStartPar
\sphinxincludegraphics{{production_deepset_training_curves}.png}

\sphinxAtStartPar
\sphinxincludegraphics{{future_step_accuracy}.png}


\subsection{Exp 03 — MTL Set Transformer}
\label{\detokenize{training:exp-03-mtl-set-transformer}}
\sphinxAtStartPar
Introduces self\sphinxhyphen{}attention over the cell dimension. RSRP regression diagnostic:

\sphinxAtStartPar
\sphinxincludegraphics{{rsrp_regression_diagnostic}.png}

\sphinxAtStartPar
\sphinxincludegraphics{{mtl_training_curves}.png}


\subsection{Exp 04 — 6G Predictive (Optuna\sphinxhyphen{}Tuned)}
\label{\detokenize{training:exp-04-6g-predictive-optuna-tuned}}
\sphinxAtStartPar
Optuna runs 50 trials with Tree\sphinxhyphen{}structured Parzen Estimator (TPE) sampler to find optimal \sphinxcode{\sphinxupquote{LSTM\_UNITS}}, \sphinxcode{\sphinxupquote{N\_HEADS}}, \sphinxcode{\sphinxupquote{DROPOUT}}, and \sphinxcode{\sphinxupquote{LR\_INIT}}. Best epoch: 19.

\sphinxAtStartPar
\sphinxincludegraphics{{6g_training_curves}.png}


\subsection{Exp 05 — Strategic DeepSet (Context Injection)}
\label{\detokenize{training:exp-05-strategic-deepset-context-injection}}
\sphinxAtStartPar
Injects global UE context (\sphinxcode{\sphinxupquote{speed}}, \sphinxcode{\sphinxupquote{dir\_cos}}, \sphinxcode{\sphinxupquote{dir\_sin}}, \sphinxcode{\sphinxupquote{cell\_load}}, \sphinxcode{\sphinxupquote{ho\_class\_onehot}}) into the cell aggregation to make the model mobility\sphinxhyphen{}regime\sphinxhyphen{}aware.

\sphinxAtStartPar
\sphinxincludegraphics{{strategic_training_curves}.png}


\subsection{Exp 06 — Ray\sphinxhyphen{}Tracing Fine\sphinxhyphen{}Tune (Progressive Unfreezing)}
\label{\detokenize{training:exp-06-ray-tracing-fine-tune-progressive-unfreezing}}
\sphinxAtStartPar
Three\sphinxhyphen{}phase progressive fine\sphinxhyphen{}tuning on the ray\sphinxhyphen{}tracing domain. Uses a temperature\sphinxhyphen{}scaled calibration (T=1.1) post\sphinxhyphen{}training.

\begin{sphinxVerbatim}[commandchars=\\\{\}]
Phase 1 (25 ep, LR=1e\PYGZhy{}4): head only   \(\rightarrow\) val top\PYGZhy{}3 ↑
Phase 2 (25 ep, LR=5e\PYGZhy{}5): upper blocks \(\rightarrow\) continued improvement
Phase 3 (15 ep, LR=1e\PYGZhy{}6): full model  \(\rightarrow\) final adaptation
\end{sphinxVerbatim}

\sphinxAtStartPar
\sphinxincludegraphics{{i_static/mage/ray_finetune_training_curves}.png}


\subsection{Exp 07 — Ray\sphinxhyphen{}ID CrossFormer (SOTA / Production)}
\label{\detokenize{training:exp-07-ray-id-crossformer-sota-production}}
\sphinxAtStartPar
The production training run. Trained directly on the ray\sphinxhyphen{}tracing cache with cell\sphinxhyphen{}ID embeddings. Converges in \textasciitilde{}60\textendash{}70 effective epochs under early stopping.

\begin{sphinxVerbatim}[commandchars=\\\{\}]
\PYG{c+c1}{\PYGZsh{} Training (from project root, inside the tda\PYGZhy{}handover conda env)}
conda\PYG{+w}{ }run\PYG{+w}{ }\PYGZhy{}n\PYG{+w}{ }tda\PYGZhy{}handover\PYG{+w}{ }python\PYG{+w}{ }\PYGZhy{}m\PYG{+w}{ }notebooks.modeling.07\PYGZus{}ray\PYGZus{}id\PYGZus{}crossformer
\end{sphinxVerbatim}

\sphinxAtStartPar
\sphinxincludegraphics{{ray_id_crossformer_training_curves}.png}

\sphinxAtStartPar
\sphinxincludegraphics{{ray_id_crossformer_confusion_matrix}.png}


\bigskip\hrule\bigskip



\section{6. Evaluation Protocol}
\label{\detokenize{training:evaluation-protocol}}

\subsection{6.1 Metrics}
\label{\detokenize{training:metrics}}

\begin{savenotes}\sphinxattablestart
\sphinxthistablewithglobalstyle
\centering
\begin{tabulary}{\linewidth}[t]{TTT}
\sphinxtoprule
\sphinxstyletheadfamily 
\sphinxAtStartPar
Metric
&\sphinxstyletheadfamily 
\sphinxAtStartPar
Formula
&\sphinxstyletheadfamily 
\sphinxAtStartPar
Role
\\
\sphinxmidrule
\sphinxtableatstartofbodyhook
\sphinxAtStartPar
\sphinxstylestrong{Top\sphinxhyphen{}1 Accuracy}
&
\sphinxAtStartPar
\(\Pr[\hat{y}_1 = y]\)
&
\sphinxAtStartPar
Exact match — hardest metric
\\
\sphinxhline
\sphinxAtStartPar
\sphinxstylestrong{Top\sphinxhyphen{}3 Accuracy}
&
\sphinxAtStartPar
\(\Pr[y \in \{\hat{y}_1, \hat{y}_2, \hat{y}_3\}]\)
&
\sphinxAtStartPar
\sphinxstylestrong{Primary production KPI}
\\
\sphinxhline
\sphinxAtStartPar
\sphinxstylestrong{Top\sphinxhyphen{}5 Accuracy}
&
\sphinxAtStartPar
\(\Pr[y \in \text{top-5}]\)
&
\sphinxAtStartPar
Coverage ceiling
\\
\sphinxhline
\sphinxAtStartPar
RSRP MAE (dBm)
&
\sphinxAtStartPar
\$\textbackslash{}text\{E\}{[}
&
\sphinxAtStartPar
\textbackslash{}hat\{r\} \sphinxhyphen{} r
\\
\sphinxbottomrule
\end{tabulary}
\sphinxtableafterendhook\par
\sphinxattableend\end{savenotes}


\subsection{6.2 Why Top\sphinxhyphen{}3 Is the Primary Production Metric}
\label{\detokenize{training:why-top-3-is-the-primary-production-metric}}
\sphinxAtStartPar
In a live 5G NR deployment, the edge\sphinxhyphen{}AI inference does not issue a single hard handover command. Instead, it returns a \sphinxstylestrong{ranked candidate list} to the RAN controller (via X2/Xn interface), which then selects the best reachable candidate based on:
\begin{itemize}
\item {} 
\sphinxAtStartPar
live X2/Xn resource availability,

\item {} 
\sphinxAtStartPar
QoS profile matching,

\item {} 
\sphinxAtStartPar
load constraints on the candidate cell.

\end{itemize}

\sphinxAtStartPar
This means the network can exploit the \sphinxstylestrong{top\sphinxhyphen{}3 predicted candidates}, not just the single best prediction. A Top\sphinxhyphen{}3 accuracy of \sphinxstylestrong{90.87 \%} (CrossFormer) means the optimal target cell is presented in the candidate shortlist in 9 out of 10 handover events, essentially eliminating blind handover failures in production.


\subsection{6.3 Evaluation Script}
\label{\detokenize{training:evaluation-script}}
\begin{sphinxVerbatim}[commandchars=\\\{\}]
\PYG{c+c1}{\PYGZsh{} Evaluate CrossFormer on ray\PYGZhy{}tracing test split}
conda\PYG{+w}{ }run\PYG{+w}{ }\PYGZhy{}n\PYG{+w}{ }tda\PYGZhy{}handover\PYG{+w}{ }python\PYG{+w}{ }src/evaluate\PYGZus{}model.py\PYG{+w}{ }\PYG{l+s+se}{\PYGZbs{}}
\PYG{+w}{  }\PYGZhy{}\PYGZhy{}model\PYG{+w}{ }models/best\PYGZus{}ray\PYGZus{}id\PYGZus{}crossformer.keras\PYG{+w}{ }\PYG{l+s+se}{\PYGZbs{}}
\PYG{+w}{  }\PYGZhy{}\PYGZhy{}label\PYG{+w}{ }optimal
\end{sphinxVerbatim}


\subsection{6.4 Per\sphinxhyphen{}Cell Recall}
\label{\detokenize{training:per-cell-recall}}
\sphinxAtStartPar
Per\sphinxhyphen{}cell recall plots expose which cells are systematically confused — useful for identifying coverage holes, highly loaded cells, or geo\sphinxhyphen{}clusters with ambiguous candidates:

\sphinxAtStartPar
\sphinxincludegraphics{{temporal_deepset_per_cell_recall}.png}


\bigskip\hrule\bigskip



\section{7. Production Source Modules}
\label{\detokenize{training:production-source-modules}}
\sphinxAtStartPar
The \sphinxcode{\sphinxupquote{src/production/}} directory contains reusable Python modules that mirror notebook logic in a script\sphinxhyphen{}friendly, reproducible form:


\begin{savenotes}\sphinxattablestart
\sphinxthistablewithglobalstyle
\centering
\begin{tabulary}{\linewidth}[t]{TT}
\sphinxtoprule
\sphinxstyletheadfamily 
\sphinxAtStartPar
Module
&\sphinxstyletheadfamily 
\sphinxAtStartPar
Description
\\
\sphinxmidrule
\sphinxtableatstartofbodyhook
\sphinxAtStartPar
\sphinxcode{\sphinxupquote{src/production/temporal\_deepset\_data.py}}
&
\sphinxAtStartPar
Windowed cache builder + scaler for DeepSet experiments
\\
\sphinxhline
\sphinxAtStartPar
\sphinxcode{\sphinxupquote{src/production/temporal\_deepset\_model.py}}
&
\sphinxAtStartPar
Keras model builder + loader (with custom \sphinxcode{\sphinxupquote{MaskedGlobalAveragePooling}})
\\
\sphinxhline
\sphinxAtStartPar
\sphinxcode{\sphinxupquote{src/production/temporal\_deepset\_decision.py}}
&
\sphinxAtStartPar
TTT + margin + cooldown stability gate
\\
\sphinxhline
\sphinxAtStartPar
\sphinxcode{\sphinxupquote{src/production/mh\_transformer\_data.py}}
&
\sphinxAtStartPar
Multi\sphinxhyphen{}horizon cache builder for transformer experiments
\\
\sphinxhline
\sphinxAtStartPar
\sphinxcode{\sphinxupquote{src/production/mh\_transformer\_model.py}}
&
\sphinxAtStartPar
Multi\sphinxhyphen{}horizon transformer model builder
\\
\sphinxhline
\sphinxAtStartPar
\sphinxcode{\sphinxupquote{src/train\_model.py}}
&
\sphinxAtStartPar
Train TemporalDeepSet (single\sphinxhyphen{}step baseline)
\\
\sphinxhline
\sphinxAtStartPar
\sphinxcode{\sphinxupquote{src/train\_mh\_transformer.py}}
&
\sphinxAtStartPar
Train multi\sphinxhyphen{}horizon transformer
\\
\sphinxhline
\sphinxAtStartPar
\sphinxcode{\sphinxupquote{src/evaluate\_model.py}}
&
\sphinxAtStartPar
Evaluate and write metrics JSON
\\
\sphinxhline
\sphinxAtStartPar
\sphinxcode{\sphinxupquote{src/inference.py}}
&
\sphinxAtStartPar
Offline streaming inference + stability gate
\\
\sphinxbottomrule
\end{tabulary}
\sphinxtableafterendhook\par
\sphinxattableend\end{savenotes}
\begin{quote}

\sphinxAtStartPar
\sphinxstylestrong{Inference entry point uses \sphinxcode{\sphinxupquote{models/best\_ray\_id\_crossformer.keras}}} as the default production model. The \sphinxcode{\sphinxupquote{src/inference.py}} script runs per\sphinxhyphen{}UE rolling\sphinxhyphen{}window inference with the stability gate (TTT=3 steps, margin=0.15, cooldown=5 steps).
\end{quote}

\sphinxstepscope


\chapter{5G\sphinxhyphen{}to\sphinxhyphen{}6G Intelligent Handover \& Cell Selection Prediction Engine}
\label{\detokenize{deployement:g-to-6g-intelligent-handover-cell-selection-prediction-engine}}\label{\detokenize{deployement::doc}}

\section{Production Deployment Plan \& Architectural Specification}
\label{\detokenize{deployement:production-deployment-plan-architectural-specification}}
\sphinxAtStartPar
\sphinxstylestrong{Document Version:} 1.0.0\\
\sphinxstylestrong{Status:} Ready for Telco Infrastructure Review\\
\sphinxstylestrong{Target Environment:} Next\sphinxhyphen{}Generation Radio Access Network (NG\sphinxhyphen{}RAN) / gNodeB Server Infrastructure\\
\sphinxstylestrong{Deployment Context:} 5G Non\sphinxhyphen{}Standalone (NSA) Anchored with 6G\sphinxhyphen{}Forward Transition Capabilities


\bigskip\hrule\bigskip



\section{1. Executive Summary \& Context}
\label{\detokenize{deployement:executive-summary-context}}
\sphinxAtStartPar
This deployment plan outlines the integration, operational architecture, and performance verification for an intelligent \sphinxstylestrong{Multi\sphinxhyphen{}Horizon Handover (HO) and Cell Selection Prediction Engine}. Operating over a cluster of \sphinxstylestrong{10 heterogeneous macro and micro cells}, this system mitigates the latency and ping\sphinxhyphen{}pong effects typical of traditional reactive handover mechanisms (A3/A5 event\sphinxhyphen{}triggered).

\sphinxAtStartPar
By leveraging deep learning models validated on a highly precise simulation framework, the engine predicts cell suitability \sphinxstylestrong{5 timestamps into the future (250ms lookahead with 50ms sampling intervals)}. This proactive approach allows telecom operators to pre\sphinxhyphen{}allocate radio resources and execute seamless context transfers, providing a foundational bridge for the transition from 5G New Radio (NR) to early 6G ultra\sphinxhyphen{}reliable low\sphinxhyphen{}latency communication (URLLC) environments.


\subsection{Current Network Implementation Matrix}
\label{\detokenize{deployement:current-network-implementation-matrix}}\begin{itemize}
\item {} 
\sphinxAtStartPar
\sphinxstylestrong{Deployment Mode:} 5G Non\sphinxhyphen{}Standalone (NSA) Architecture (Option 3x). Control plane anchored via 4G Evolved Packet Core (EPC / eNodeB), user plane boosted via 5G NR servers (gNodeB).

\item {} 
\sphinxAtStartPar
\sphinxstylestrong{Target Scale:} 10 highly dynamic cells, including aerial nodes.

\item {} 
\sphinxAtStartPar
\sphinxstylestrong{Core Application:} Proactive Conditional Handover (CHO) optimization and load\sphinxhyphen{}aware cell selection.

\end{itemize}


\bigskip\hrule\bigskip



\section{2. Advanced Simulation \& Fine\sphinxhyphen{}Tuning Pipeline}
\label{\detokenize{deployement:advanced-simulation-fine-tuning-pipeline}}
\sphinxAtStartPar
The underlying model was developed, validated, and fine\sphinxhyphen{}tuned using a high\sphinxhyphen{}fidelity hybrid simulation pipeline to mirror realistic radio propagation and subscriber mobility:

\begin{sphinxVerbatim}[commandchars=\\\{\}]
\PYG{o}{+}\PYG{o}{\PYGZhy{}}\PYG{o}{\PYGZhy{}}\PYG{o}{\PYGZhy{}}\PYG{o}{\PYGZhy{}}\PYG{o}{\PYGZhy{}}\PYG{o}{\PYGZhy{}}\PYG{o}{\PYGZhy{}}\PYG{o}{\PYGZhy{}}\PYG{o}{\PYGZhy{}}\PYG{o}{\PYGZhy{}}\PYG{o}{\PYGZhy{}}\PYG{o}{\PYGZhy{}}\PYG{o}{\PYGZhy{}}\PYG{o}{\PYGZhy{}}\PYG{o}{\PYGZhy{}}\PYG{o}{\PYGZhy{}}\PYG{o}{\PYGZhy{}}\PYG{o}{\PYGZhy{}}\PYG{o}{\PYGZhy{}}\PYG{o}{\PYGZhy{}}\PYG{o}{\PYGZhy{}}\PYG{o}{\PYGZhy{}}\PYG{o}{\PYGZhy{}}\PYG{o}{\PYGZhy{}}\PYG{o}{\PYGZhy{}}\PYG{o}{\PYGZhy{}}\PYG{o}{\PYGZhy{}}\PYG{o}{\PYGZhy{}}\PYG{o}{\PYGZhy{}}\PYG{o}{\PYGZhy{}}\PYG{o}{\PYGZhy{}}\PYG{o}{\PYGZhy{}}\PYG{o}{\PYGZhy{}}\PYG{o}{\PYGZhy{}}\PYG{o}{\PYGZhy{}}\PYG{o}{\PYGZhy{}}\PYG{o}{\PYGZhy{}}\PYG{o}{\PYGZhy{}}\PYG{o}{\PYGZhy{}}\PYG{o}{+}
\PYG{o}{|}  \PYG{n}{ns}\PYG{o}{\PYGZhy{}}\PYG{l+m+mi}{3} \PYG{n}{Network} \PYG{n}{Simulation} \PYG{n}{Framework}    \PYG{o}{|} \PYG{o}{\PYGZhy{}}\PYG{o}{\PYGZhy{}}\PYG{o}{\PYGZhy{}}\PYG{o}{\PYGZgt{}} \PYG{n}{Microscopic} \PYG{n}{Mobility} \PYG{n}{Models}
\PYG{o}{|}  \PYG{p}{(}\PYG{n}{Macroscopic} \PYG{n}{Network} \PYG{o}{\PYGZam{}} \PYG{n}{Topology}\PYG{p}{)}     \PYG{o}{|}      \PYG{o}{\PYGZam{}} \PYG{n}{User} \PYG{n}{Equipment} \PYG{p}{(}\PYG{n}{UE}\PYG{p}{)} \PYG{n}{Trajectories}
\PYG{o}{+}\PYG{o}{\PYGZhy{}}\PYG{o}{\PYGZhy{}}\PYG{o}{\PYGZhy{}}\PYG{o}{\PYGZhy{}}\PYG{o}{\PYGZhy{}}\PYG{o}{\PYGZhy{}}\PYG{o}{\PYGZhy{}}\PYG{o}{\PYGZhy{}}\PYG{o}{\PYGZhy{}}\PYG{o}{\PYGZhy{}}\PYG{o}{\PYGZhy{}}\PYG{o}{\PYGZhy{}}\PYG{o}{\PYGZhy{}}\PYG{o}{\PYGZhy{}}\PYG{o}{\PYGZhy{}}\PYG{o}{\PYGZhy{}}\PYG{o}{\PYGZhy{}}\PYG{o}{\PYGZhy{}}\PYG{o}{\PYGZhy{}}\PYG{o}{\PYGZhy{}}\PYG{o}{\PYGZhy{}}\PYG{o}{\PYGZhy{}}\PYG{o}{\PYGZhy{}}\PYG{o}{\PYGZhy{}}\PYG{o}{\PYGZhy{}}\PYG{o}{\PYGZhy{}}\PYG{o}{\PYGZhy{}}\PYG{o}{\PYGZhy{}}\PYG{o}{\PYGZhy{}}\PYG{o}{\PYGZhy{}}\PYG{o}{\PYGZhy{}}\PYG{o}{\PYGZhy{}}\PYG{o}{\PYGZhy{}}\PYG{o}{\PYGZhy{}}\PYG{o}{\PYGZhy{}}\PYG{o}{\PYGZhy{}}\PYG{o}{\PYGZhy{}}\PYG{o}{\PYGZhy{}}\PYG{o}{\PYGZhy{}}\PYG{o}{+}
                    \PYG{o}{|}
                    \PYG{n}{v}
\PYG{o}{+}\PYG{o}{\PYGZhy{}}\PYG{o}{\PYGZhy{}}\PYG{o}{\PYGZhy{}}\PYG{o}{\PYGZhy{}}\PYG{o}{\PYGZhy{}}\PYG{o}{\PYGZhy{}}\PYG{o}{\PYGZhy{}}\PYG{o}{\PYGZhy{}}\PYG{o}{\PYGZhy{}}\PYG{o}{\PYGZhy{}}\PYG{o}{\PYGZhy{}}\PYG{o}{\PYGZhy{}}\PYG{o}{\PYGZhy{}}\PYG{o}{\PYGZhy{}}\PYG{o}{\PYGZhy{}}\PYG{o}{\PYGZhy{}}\PYG{o}{\PYGZhy{}}\PYG{o}{\PYGZhy{}}\PYG{o}{\PYGZhy{}}\PYG{o}{\PYGZhy{}}\PYG{o}{\PYGZhy{}}\PYG{o}{\PYGZhy{}}\PYG{o}{\PYGZhy{}}\PYG{o}{\PYGZhy{}}\PYG{o}{\PYGZhy{}}\PYG{o}{\PYGZhy{}}\PYG{o}{\PYGZhy{}}\PYG{o}{\PYGZhy{}}\PYG{o}{\PYGZhy{}}\PYG{o}{\PYGZhy{}}\PYG{o}{\PYGZhy{}}\PYG{o}{\PYGZhy{}}\PYG{o}{\PYGZhy{}}\PYG{o}{\PYGZhy{}}\PYG{o}{\PYGZhy{}}\PYG{o}{\PYGZhy{}}\PYG{o}{\PYGZhy{}}\PYG{o}{\PYGZhy{}}\PYG{o}{\PYGZhy{}}\PYG{o}{+}
\PYG{o}{|}   \PYG{n}{Sionna} \PYG{n}{Link}\PYG{o}{\PYGZhy{}}\PYG{n}{Level} \PYG{n}{Simulator}         \PYG{o}{|} \PYG{o}{\PYGZhy{}}\PYG{o}{\PYGZhy{}}\PYG{o}{\PYGZhy{}}\PYG{o}{\PYGZgt{}} \PYG{n}{Differentiable} \PYG{n}{Physical} \PYG{n}{Layer} \PYG{p}{(}\PYG{n}{PHY}\PYG{p}{)}
\PYG{o}{|}   \PYG{p}{(}\PYG{n}{Python}\PYG{o}{\PYGZhy{}}\PYG{n}{based} \PYG{n}{Tensor} \PYG{n}{Framework}\PYG{p}{)}     \PYG{o}{|}      \PYG{n}{Urban} \PYG{n}{Micro} \PYG{p}{(}\PYG{n}{UMi}\PYG{p}{)} \PYG{n}{Channel} \PYG{n}{Profiles}
\PYG{o}{+}\PYG{o}{\PYGZhy{}}\PYG{o}{\PYGZhy{}}\PYG{o}{\PYGZhy{}}\PYG{o}{\PYGZhy{}}\PYG{o}{\PYGZhy{}}\PYG{o}{\PYGZhy{}}\PYG{o}{\PYGZhy{}}\PYG{o}{\PYGZhy{}}\PYG{o}{\PYGZhy{}}\PYG{o}{\PYGZhy{}}\PYG{o}{\PYGZhy{}}\PYG{o}{\PYGZhy{}}\PYG{o}{\PYGZhy{}}\PYG{o}{\PYGZhy{}}\PYG{o}{\PYGZhy{}}\PYG{o}{\PYGZhy{}}\PYG{o}{\PYGZhy{}}\PYG{o}{\PYGZhy{}}\PYG{o}{\PYGZhy{}}\PYG{o}{\PYGZhy{}}\PYG{o}{\PYGZhy{}}\PYG{o}{\PYGZhy{}}\PYG{o}{\PYGZhy{}}\PYG{o}{\PYGZhy{}}\PYG{o}{\PYGZhy{}}\PYG{o}{\PYGZhy{}}\PYG{o}{\PYGZhy{}}\PYG{o}{\PYGZhy{}}\PYG{o}{\PYGZhy{}}\PYG{o}{\PYGZhy{}}\PYG{o}{\PYGZhy{}}\PYG{o}{\PYGZhy{}}\PYG{o}{\PYGZhy{}}\PYG{o}{\PYGZhy{}}\PYG{o}{\PYGZhy{}}\PYG{o}{\PYGZhy{}}\PYG{o}{\PYGZhy{}}\PYG{o}{\PYGZhy{}}\PYG{o}{\PYGZhy{}}\PYG{o}{+}
                    \PYG{o}{|}
                    \PYG{n}{v} \PYG{p}{(}\PYG{n}{Fine}\PYG{o}{\PYGZhy{}}\PYG{n}{Tuning} \PYG{n}{Layer}\PYG{p}{)}
\PYG{o}{+}\PYG{o}{\PYGZhy{}}\PYG{o}{\PYGZhy{}}\PYG{o}{\PYGZhy{}}\PYG{o}{\PYGZhy{}}\PYG{o}{\PYGZhy{}}\PYG{o}{\PYGZhy{}}\PYG{o}{\PYGZhy{}}\PYG{o}{\PYGZhy{}}\PYG{o}{\PYGZhy{}}\PYG{o}{\PYGZhy{}}\PYG{o}{\PYGZhy{}}\PYG{o}{\PYGZhy{}}\PYG{o}{\PYGZhy{}}\PYG{o}{\PYGZhy{}}\PYG{o}{\PYGZhy{}}\PYG{o}{\PYGZhy{}}\PYG{o}{\PYGZhy{}}\PYG{o}{\PYGZhy{}}\PYG{o}{\PYGZhy{}}\PYG{o}{\PYGZhy{}}\PYG{o}{\PYGZhy{}}\PYG{o}{\PYGZhy{}}\PYG{o}{\PYGZhy{}}\PYG{o}{\PYGZhy{}}\PYG{o}{\PYGZhy{}}\PYG{o}{\PYGZhy{}}\PYG{o}{\PYGZhy{}}\PYG{o}{\PYGZhy{}}\PYG{o}{\PYGZhy{}}\PYG{o}{\PYGZhy{}}\PYG{o}{\PYGZhy{}}\PYG{o}{\PYGZhy{}}\PYG{o}{\PYGZhy{}}\PYG{o}{\PYGZhy{}}\PYG{o}{\PYGZhy{}}\PYG{o}{\PYGZhy{}}\PYG{o}{\PYGZhy{}}\PYG{o}{\PYGZhy{}}\PYG{o}{\PYGZhy{}}\PYG{o}{+}
\PYG{o}{|}       \PYG{n}{Blender} \PYG{l+m+mi}{3}\PYG{n}{D} \PYG{n}{Environment}          \PYG{o}{|} \PYG{o}{\PYGZhy{}}\PYG{o}{\PYGZhy{}}\PYG{o}{\PYGZhy{}}\PYG{o}{\PYGZgt{}} \PYG{n}{Accurate} \PYG{n}{Urban} \PYG{n}{Topology} \PYG{n}{Generation}
\PYG{o}{|}   \PYG{p}{(}\PYG{n}{CFC} \PYG{n}{Mesh} \PYG{n}{Model} \PYG{n}{Reconstruction}\PYG{p}{)}    \PYG{o}{|}      \PYG{o}{\PYGZam{}} \PYG{n}{XML} \PYG{n}{Ray}\PYG{o}{\PYGZhy{}}\PYG{n}{Tracing} \PYG{n}{Map} \PYG{n}{Exports}
\PYG{o}{+}\PYG{o}{\PYGZhy{}}\PYG{o}{\PYGZhy{}}\PYG{o}{\PYGZhy{}}\PYG{o}{\PYGZhy{}}\PYG{o}{\PYGZhy{}}\PYG{o}{\PYGZhy{}}\PYG{o}{\PYGZhy{}}\PYG{o}{\PYGZhy{}}\PYG{o}{\PYGZhy{}}\PYG{o}{\PYGZhy{}}\PYG{o}{\PYGZhy{}}\PYG{o}{\PYGZhy{}}\PYG{o}{\PYGZhy{}}\PYG{o}{\PYGZhy{}}\PYG{o}{\PYGZhy{}}\PYG{o}{\PYGZhy{}}\PYG{o}{\PYGZhy{}}\PYG{o}{\PYGZhy{}}\PYG{o}{\PYGZhy{}}\PYG{o}{\PYGZhy{}}\PYG{o}{\PYGZhy{}}\PYG{o}{\PYGZhy{}}\PYG{o}{\PYGZhy{}}\PYG{o}{\PYGZhy{}}\PYG{o}{\PYGZhy{}}\PYG{o}{\PYGZhy{}}\PYG{o}{\PYGZhy{}}\PYG{o}{\PYGZhy{}}\PYG{o}{\PYGZhy{}}\PYG{o}{\PYGZhy{}}\PYG{o}{\PYGZhy{}}\PYG{o}{\PYGZhy{}}\PYG{o}{\PYGZhy{}}\PYG{o}{\PYGZhy{}}\PYG{o}{\PYGZhy{}}\PYG{o}{\PYGZhy{}}\PYG{o}{\PYGZhy{}}\PYG{o}{\PYGZhy{}}\PYG{o}{\PYGZhy{}}\PYG{o}{+}
\end{sphinxVerbatim}
\begin{enumerate}
\sphinxsetlistlabels{\arabic}{enumi}{enumii}{}{.}%
\item {} 
\sphinxAtStartPar
\sphinxstylestrong{Macroscopic Mobility \& Protocol Stack (\sphinxcode{\sphinxupquote{ns\sphinxhyphen{}3}}):} Packet\sphinxhyphen{}level interactions, user equipment (UE) trajectories, and macro network behavior were simulated using \sphinxcode{\sphinxupquote{ns\sphinxhyphen{}3}}.

\item {} 
\sphinxAtStartPar
\sphinxstylestrong{Physical Layer (PHY) Processing (\sphinxcode{\sphinxupquote{Sionna}}):} The python\sphinxhyphen{}based Sionna framework executed differentiable link\sphinxhyphen{}level simulations using standard 3GPP \sphinxstylestrong{Urban Micro (UMi)} channel models.

\item {} 
\sphinxAtStartPar
\sphinxstylestrong{Drone Integration (Load\sphinxhyphen{}Aware Strategy):} Unmanned Aerial Vehicles (UAVs) acting as airborne base stations were embedded into the network. A specialized load\sphinxhyphen{}aware strategy dynamically adjusted cell selection boundaries based on ground\sphinxhyphen{}user density and drone battery/backhaul constraints.

\item {} 
\sphinxAtStartPar
\sphinxstylestrong{Site\sphinxhyphen{}Specific Ray\sphinxhyphen{}Tracing Fine\sphinxhyphen{}Tuning:} To transition the model from idealized geometry to site\sphinxhyphen{}specific reality, precise \sphinxstylestrong{Computational Fluid Dynamics (CFC) meshes} of urban environments were modeled in Blender. These were exported as structured XML scene descriptions into Sionna’s ray\sphinxhyphen{}tracing engine, producing realistic multipath components, shadow fading, and diffraction profiles.

\end{enumerate}


\bigskip\hrule\bigskip



\section{3. Performance Metrics \& Telecom Significance}
\label{\detokenize{deployement:performance-metrics-telecom-significance}}
\sphinxAtStartPar
The predictive engine shifts the operational paradigm from single\sphinxhyphen{}step reactive assessment to multi\sphinxhyphen{}horizon predictive scoring. The model’s evaluation metrics demonstrate high reliability across complex scenarios:


\subsection{Performance Comparison Matrix}
\label{\detokenize{deployement:performance-comparison-matrix}}

\begin{savenotes}\sphinxattablestart
\sphinxthistablewithglobalstyle
\centering
\begin{tabulary}{\linewidth}[t]{TTTTT}
\sphinxtoprule
\sphinxstyletheadfamily 
\sphinxAtStartPar
Topology / Horizon Profile
&\sphinxstyletheadfamily 
\sphinxAtStartPar
Lookahead Interval
&\sphinxstyletheadfamily 
\sphinxAtStartPar
Top\sphinxhyphen{}1 Accuracy Score
&\sphinxstyletheadfamily 
\sphinxAtStartPar
Top\sphinxhyphen{}3 Accuracy Score
&\sphinxstyletheadfamily 
\sphinxAtStartPar
Operational Application
\\
\sphinxmidrule
\sphinxtableatstartofbodyhook
\sphinxAtStartPar
\sphinxstylestrong{Perfect Hexagonal Topology} (Ideal Baseline)
&
\sphinxAtStartPar
1\sphinxhyphen{}Step (50ms)
&
\sphinxAtStartPar
\textendash{}
&
\sphinxAtStartPar
\sphinxstylestrong{98.0\%}
&
\sphinxAtStartPar
Pure theoretical optimization / static verification.
\\
\sphinxhline
\sphinxAtStartPar
\sphinxstylestrong{Realistic Urban Micro (UMi) + Drones} (Production)
&
\sphinxAtStartPar
5\sphinxhyphen{}Step Multi\sphinxhyphen{}Horizon (250ms)
&
\sphinxAtStartPar
\sphinxstylestrong{56.0\%}
&
\sphinxAtStartPar
\sphinxstylestrong{85.0\%}
&
\sphinxAtStartPar
\sphinxstylestrong{Active Production Configuration} for Conditional Handover.
\\
\sphinxbottomrule
\end{tabulary}
\sphinxtableafterendhook\par
\sphinxattableend\end{savenotes}


\subsection{The Telecommunication Benefit of Top\sphinxhyphen{}3 Accuracy (85\%)}
\label{\detokenize{deployement:the-telecommunication-benefit-of-top-3-accuracy-85}}
\sphinxAtStartPar
While a \sphinxstylestrong{56\% Top\sphinxhyphen{}1 accuracy} at a 250ms horizon seems low for standard classification tasks, in cellular networking infrastructure, the \sphinxstylestrong{85\% Top\sphinxhyphen{}3 accuracy is a breakthrough operational metric}.
\begin{itemize}
\item {} 
\sphinxAtStartPar
\sphinxstylestrong{Why Top\sphinxhyphen{}1 Suffers:} At 250ms in the future within highly dynamic UMi environments (especially with moving drones and ray\sphinxhyphen{}traced shadow fading), fast\sphinxhyphen{}fading effects and blockages make predicting the \sphinxstyleemphasis{absolute absolute best} single cell highly volatile.

\item {} 
\sphinxAtStartPar
\sphinxstylestrong{The Power of Top\sphinxhyphen{}3 Execution:} Telecom standards (3GPP Release 16/17) support \sphinxstylestrong{Conditional Handover (CHO)}. Instead of preparing a single target cell, the gNodeB can pre\sphinxhyphen{}prepare a candidate cell pool.

\item {} 
\sphinxAtStartPar
\sphinxstylestrong{Resource Optimization:} Achieving 85\% Top\sphinxhyphen{}3 accuracy means that 85\% of the time, the optimal target cell is within our prepared pool of 3 cells. The NR server can pre\sphinxhyphen{}allocate random\sphinxhyphen{}access channel (RACH) resources and initiate early context transfers for these 3 cells. This eliminates the latency of fetching user context \sphinxstyleemphasis{after} a link degradation event, reducing Handover Failure (HOF) rates and Radio Link Failures (RLF) to near zero without overloading the backhaul with excessive candidate preparations.

\end{itemize}


\bigskip\hrule\bigskip



\section{4. Target Deployment Architecture}
\label{\detokenize{deployement:target-deployment-architecture}}
\sphinxAtStartPar
The prediction engine runs containerized at the edge of the Next\sphinxhyphen{}Generation Radio Access Network (NG\sphinxhyphen{}RAN), interacting directly with the RRC (Radio Resource Control) layer of the NR servers.

\begin{sphinxVerbatim}[commandchars=\\\{\}]
\PYG{o}{+}\PYG{o}{\PYGZhy{}}\PYG{o}{\PYGZhy{}}\PYG{o}{\PYGZhy{}}\PYG{o}{\PYGZhy{}}\PYG{o}{\PYGZhy{}}\PYG{o}{\PYGZhy{}}\PYG{o}{\PYGZhy{}}\PYG{o}{\PYGZhy{}}\PYG{o}{\PYGZhy{}}\PYG{o}{\PYGZhy{}}\PYG{o}{\PYGZhy{}}\PYG{o}{\PYGZhy{}}\PYG{o}{\PYGZhy{}}\PYG{o}{\PYGZhy{}}\PYG{o}{\PYGZhy{}}\PYG{o}{\PYGZhy{}}\PYG{o}{\PYGZhy{}}\PYG{o}{\PYGZhy{}}\PYG{o}{\PYGZhy{}}\PYG{o}{\PYGZhy{}}\PYG{o}{\PYGZhy{}}\PYG{o}{\PYGZhy{}}\PYG{o}{\PYGZhy{}}\PYG{o}{\PYGZhy{}}\PYG{o}{\PYGZhy{}}\PYG{o}{\PYGZhy{}}\PYG{o}{\PYGZhy{}}\PYG{o}{\PYGZhy{}}\PYG{o}{\PYGZhy{}}\PYG{o}{\PYGZhy{}}\PYG{o}{\PYGZhy{}}\PYG{o}{\PYGZhy{}}\PYG{o}{\PYGZhy{}}\PYG{o}{\PYGZhy{}}\PYG{o}{\PYGZhy{}}\PYG{o}{\PYGZhy{}}\PYG{o}{\PYGZhy{}}\PYG{o}{\PYGZhy{}}\PYG{o}{\PYGZhy{}}\PYG{o}{\PYGZhy{}}\PYG{o}{\PYGZhy{}}\PYG{o}{\PYGZhy{}}\PYG{o}{\PYGZhy{}}\PYG{o}{\PYGZhy{}}\PYG{o}{\PYGZhy{}}\PYG{o}{\PYGZhy{}}\PYG{o}{\PYGZhy{}}\PYG{o}{\PYGZhy{}}\PYG{o}{\PYGZhy{}}\PYG{o}{\PYGZhy{}}\PYG{o}{\PYGZhy{}}\PYG{o}{\PYGZhy{}}\PYG{o}{\PYGZhy{}}\PYG{o}{\PYGZhy{}}\PYG{o}{\PYGZhy{}}\PYG{o}{\PYGZhy{}}\PYG{o}{\PYGZhy{}}\PYG{o}{\PYGZhy{}}\PYG{o}{\PYGZhy{}}\PYG{o}{\PYGZhy{}}\PYG{o}{\PYGZhy{}}\PYG{o}{\PYGZhy{}}\PYG{o}{\PYGZhy{}}\PYG{o}{\PYGZhy{}}\PYG{o}{\PYGZhy{}}\PYG{o}{\PYGZhy{}}\PYG{o}{\PYGZhy{}}\PYG{o}{\PYGZhy{}}\PYG{o}{\PYGZhy{}}\PYG{o}{\PYGZhy{}}\PYG{o}{\PYGZhy{}}\PYG{o}{\PYGZhy{}}\PYG{o}{\PYGZhy{}}\PYG{o}{\PYGZhy{}}\PYG{o}{\PYGZhy{}}\PYG{o}{\PYGZhy{}}\PYG{o}{\PYGZhy{}}\PYG{o}{\PYGZhy{}}\PYG{o}{\PYGZhy{}}\PYG{o}{\PYGZhy{}}\PYG{o}{\PYGZhy{}}\PYG{o}{+}
\PYG{o}{|}                                \PYG{l+m+mi}{5}\PYG{n}{G} \PYG{n}{NR} \PYG{n}{SERVER} \PYG{p}{(}\PYG{n}{gNodeB}\PYG{p}{)}                            \PYG{o}{|}
\PYG{o}{|}                                                                                 \PYG{o}{|}
\PYG{o}{|}  \PYG{o}{+}\PYG{o}{\PYGZhy{}}\PYG{o}{\PYGZhy{}}\PYG{o}{\PYGZhy{}}\PYG{o}{\PYGZhy{}}\PYG{o}{\PYGZhy{}}\PYG{o}{\PYGZhy{}}\PYG{o}{\PYGZhy{}}\PYG{o}{\PYGZhy{}}\PYG{o}{\PYGZhy{}}\PYG{o}{\PYGZhy{}}\PYG{o}{\PYGZhy{}}\PYG{o}{\PYGZhy{}}\PYG{o}{\PYGZhy{}}\PYG{o}{\PYGZhy{}}\PYG{o}{\PYGZhy{}}\PYG{o}{\PYGZhy{}}\PYG{o}{\PYGZhy{}}\PYG{o}{\PYGZhy{}}\PYG{o}{\PYGZhy{}}\PYG{o}{\PYGZhy{}}\PYG{o}{\PYGZhy{}}\PYG{o}{\PYGZhy{}}\PYG{o}{\PYGZhy{}}\PYG{o}{\PYGZhy{}}\PYG{o}{\PYGZhy{}}\PYG{o}{\PYGZhy{}}\PYG{o}{\PYGZhy{}}\PYG{o}{+}                 \PYG{o}{+}\PYG{o}{\PYGZhy{}}\PYG{o}{\PYGZhy{}}\PYG{o}{\PYGZhy{}}\PYG{o}{\PYGZhy{}}\PYG{o}{\PYGZhy{}}\PYG{o}{\PYGZhy{}}\PYG{o}{\PYGZhy{}}\PYG{o}{\PYGZhy{}}\PYG{o}{\PYGZhy{}}\PYG{o}{\PYGZhy{}}\PYG{o}{\PYGZhy{}}\PYG{o}{\PYGZhy{}}\PYG{o}{\PYGZhy{}}\PYG{o}{\PYGZhy{}}\PYG{o}{\PYGZhy{}}\PYG{o}{\PYGZhy{}}\PYG{o}{\PYGZhy{}}\PYG{o}{\PYGZhy{}}\PYG{o}{\PYGZhy{}}\PYG{o}{\PYGZhy{}}\PYG{o}{\PYGZhy{}}\PYG{o}{\PYGZhy{}}\PYG{o}{\PYGZhy{}}\PYG{o}{\PYGZhy{}}\PYG{o}{\PYGZhy{}}\PYG{o}{\PYGZhy{}}\PYG{o}{\PYGZhy{}}\PYG{o}{\PYGZhy{}}\PYG{o}{\PYGZhy{}}\PYG{o}{+}  \PYG{o}{|}
\PYG{o}{|}  \PYG{o}{|}     \PYG{n}{CU}\PYG{o}{\PYGZhy{}}\PYG{n}{CP} \PYG{o}{/} \PYG{n}{RRC} \PYG{n}{Layer}     \PYG{o}{|}  \PYG{n}{UE} \PYG{n}{Metrics}     \PYG{o}{|}  \PYG{n}{Handover} \PYG{n}{Prediction} \PYG{n}{Engine} \PYG{o}{|}  \PYG{o}{|}
\PYG{o}{|}  \PYG{o}{|} \PYG{p}{(}\PYG{n}{Radio} \PYG{n}{Resource} \PYG{n}{Control}\PYG{p}{)}  \PYG{o}{|}\PYG{o}{\PYGZhy{}}\PYG{o}{\PYGZhy{}}\PYG{o}{\PYGZhy{}}\PYG{o}{\PYGZhy{}}\PYG{o}{\PYGZhy{}}\PYG{o}{\PYGZhy{}}\PYG{o}{\PYGZhy{}}\PYG{o}{\PYGZhy{}}\PYG{o}{\PYGZhy{}}\PYG{o}{\PYGZhy{}}\PYG{o}{\PYGZhy{}}\PYG{o}{\PYGZhy{}}\PYG{o}{\PYGZhy{}}\PYG{o}{\PYGZhy{}}\PYG{o}{\PYGZhy{}}\PYG{o}{\PYGZhy{}}\PYG{o}{\PYGZgt{}}\PYG{o}{|}     \PYG{p}{(}\PYG{n}{ONNX} \PYG{n}{Runtime} \PYG{n}{Container}\PYG{p}{)}\PYG{o}{|}  \PYG{o}{|}
\PYG{o}{|}  \PYG{o}{+}\PYG{o}{\PYGZhy{}}\PYG{o}{\PYGZhy{}}\PYG{o}{\PYGZhy{}}\PYG{o}{\PYGZhy{}}\PYG{o}{\PYGZhy{}}\PYG{o}{\PYGZhy{}}\PYG{o}{\PYGZhy{}}\PYG{o}{\PYGZhy{}}\PYG{o}{\PYGZhy{}}\PYG{o}{\PYGZhy{}}\PYG{o}{\PYGZhy{}}\PYG{o}{\PYGZhy{}}\PYG{o}{\PYGZhy{}}\PYG{o}{\PYGZhy{}}\PYG{o}{\PYGZhy{}}\PYG{o}{\PYGZhy{}}\PYG{o}{\PYGZhy{}}\PYG{o}{\PYGZhy{}}\PYG{o}{\PYGZhy{}}\PYG{o}{\PYGZhy{}}\PYG{o}{\PYGZhy{}}\PYG{o}{\PYGZhy{}}\PYG{o}{\PYGZhy{}}\PYG{o}{\PYGZhy{}}\PYG{o}{\PYGZhy{}}\PYG{o}{\PYGZhy{}}\PYG{o}{\PYGZhy{}}\PYG{o}{+} \PYG{p}{(}\PYG{n}{RSRP}\PYG{o}{/}\PYG{n}{RSRQ}\PYG{o}{/}\PYG{n}{CQI}\PYG{p}{)} \PYG{o}{+}\PYG{o}{\PYGZhy{}}\PYG{o}{\PYGZhy{}}\PYG{o}{\PYGZhy{}}\PYG{o}{\PYGZhy{}}\PYG{o}{\PYGZhy{}}\PYG{o}{\PYGZhy{}}\PYG{o}{\PYGZhy{}}\PYG{o}{\PYGZhy{}}\PYG{o}{\PYGZhy{}}\PYG{o}{\PYGZhy{}}\PYG{o}{\PYGZhy{}}\PYG{o}{\PYGZhy{}}\PYG{o}{\PYGZhy{}}\PYG{o}{\PYGZhy{}}\PYG{o}{\PYGZhy{}}\PYG{o}{\PYGZhy{}}\PYG{o}{\PYGZhy{}}\PYG{o}{\PYGZhy{}}\PYG{o}{\PYGZhy{}}\PYG{o}{\PYGZhy{}}\PYG{o}{\PYGZhy{}}\PYG{o}{\PYGZhy{}}\PYG{o}{\PYGZhy{}}\PYG{o}{\PYGZhy{}}\PYG{o}{\PYGZhy{}}\PYG{o}{\PYGZhy{}}\PYG{o}{\PYGZhy{}}\PYG{o}{\PYGZhy{}}\PYG{o}{\PYGZhy{}}\PYG{o}{+}  \PYG{o}{|}
\PYG{o}{|}                \PYG{o}{\PYGZca{}}                                              \PYG{o}{|}                 \PYG{o}{|}
\PYG{o}{|}                \PYG{o}{|} \PYG{n}{Execute} \PYG{n}{Proactive} \PYG{n}{Pool} \PYG{n}{Preparation}          \PYG{o}{|}                 \PYG{o}{|}
\PYG{o}{|}                \PYG{o}{+}\PYG{o}{\PYGZhy{}}\PYG{o}{\PYGZhy{}}\PYG{o}{\PYGZhy{}}\PYG{o}{\PYGZhy{}}\PYG{o}{\PYGZhy{}}\PYG{o}{\PYGZhy{}}\PYG{o}{\PYGZhy{}}\PYG{o}{\PYGZhy{}}\PYG{o}{\PYGZhy{}}\PYG{o}{\PYGZhy{}}\PYG{o}{\PYGZhy{}}\PYG{o}{\PYGZhy{}}\PYG{o}{\PYGZhy{}}\PYG{o}{\PYGZhy{}}\PYG{o}{\PYGZhy{}}\PYG{o}{\PYGZhy{}}\PYG{o}{\PYGZhy{}}\PYG{o}{\PYGZhy{}}\PYG{o}{\PYGZhy{}}\PYG{o}{\PYGZhy{}}\PYG{o}{\PYGZhy{}}\PYG{o}{\PYGZhy{}}\PYG{o}{\PYGZhy{}}\PYG{o}{\PYGZhy{}}\PYG{o}{\PYGZhy{}}\PYG{o}{\PYGZhy{}}\PYG{o}{\PYGZhy{}}\PYG{o}{\PYGZhy{}}\PYG{o}{\PYGZhy{}}\PYG{o}{\PYGZhy{}}\PYG{o}{\PYGZhy{}}\PYG{o}{\PYGZhy{}}\PYG{o}{\PYGZhy{}}\PYG{o}{\PYGZhy{}}\PYG{o}{\PYGZhy{}}\PYG{o}{\PYGZhy{}}\PYG{o}{\PYGZhy{}}\PYG{o}{\PYGZhy{}}\PYG{o}{\PYGZhy{}}\PYG{o}{\PYGZhy{}}\PYG{o}{\PYGZhy{}}\PYG{o}{\PYGZhy{}}\PYG{o}{\PYGZhy{}}\PYG{o}{\PYGZhy{}}\PYG{o}{\PYGZhy{}}\PYG{o}{\PYGZhy{}}\PYG{o}{+}                 \PYG{o}{|}
\PYG{o}{|}                               \PYG{p}{(}\PYG{n}{Top}\PYG{o}{\PYGZhy{}}\PYG{l+m+mi}{3} \PYG{n}{Cell} \PYG{n}{IDs}\PYG{p}{)}                                  \PYG{o}{|}
\PYG{o}{+}\PYG{o}{\PYGZhy{}}\PYG{o}{\PYGZhy{}}\PYG{o}{\PYGZhy{}}\PYG{o}{\PYGZhy{}}\PYG{o}{\PYGZhy{}}\PYG{o}{\PYGZhy{}}\PYG{o}{\PYGZhy{}}\PYG{o}{\PYGZhy{}}\PYG{o}{\PYGZhy{}}\PYG{o}{\PYGZhy{}}\PYG{o}{\PYGZhy{}}\PYG{o}{\PYGZhy{}}\PYG{o}{\PYGZhy{}}\PYG{o}{\PYGZhy{}}\PYG{o}{\PYGZhy{}}\PYG{o}{\PYGZhy{}}\PYG{o}{\PYGZhy{}}\PYG{o}{\PYGZhy{}}\PYG{o}{\PYGZhy{}}\PYG{o}{\PYGZhy{}}\PYG{o}{\PYGZhy{}}\PYG{o}{\PYGZhy{}}\PYG{o}{\PYGZhy{}}\PYG{o}{\PYGZhy{}}\PYG{o}{\PYGZhy{}}\PYG{o}{\PYGZhy{}}\PYG{o}{\PYGZhy{}}\PYG{o}{\PYGZhy{}}\PYG{o}{\PYGZhy{}}\PYG{o}{\PYGZhy{}}\PYG{o}{\PYGZhy{}}\PYG{o}{\PYGZhy{}}\PYG{o}{\PYGZhy{}}\PYG{o}{\PYGZhy{}}\PYG{o}{\PYGZhy{}}\PYG{o}{\PYGZhy{}}\PYG{o}{\PYGZhy{}}\PYG{o}{\PYGZhy{}}\PYG{o}{\PYGZhy{}}\PYG{o}{\PYGZhy{}}\PYG{o}{\PYGZhy{}}\PYG{o}{\PYGZhy{}}\PYG{o}{\PYGZhy{}}\PYG{o}{\PYGZhy{}}\PYG{o}{\PYGZhy{}}\PYG{o}{\PYGZhy{}}\PYG{o}{\PYGZhy{}}\PYG{o}{\PYGZhy{}}\PYG{o}{\PYGZhy{}}\PYG{o}{\PYGZhy{}}\PYG{o}{\PYGZhy{}}\PYG{o}{\PYGZhy{}}\PYG{o}{\PYGZhy{}}\PYG{o}{\PYGZhy{}}\PYG{o}{\PYGZhy{}}\PYG{o}{\PYGZhy{}}\PYG{o}{\PYGZhy{}}\PYG{o}{\PYGZhy{}}\PYG{o}{\PYGZhy{}}\PYG{o}{\PYGZhy{}}\PYG{o}{\PYGZhy{}}\PYG{o}{\PYGZhy{}}\PYG{o}{\PYGZhy{}}\PYG{o}{\PYGZhy{}}\PYG{o}{\PYGZhy{}}\PYG{o}{\PYGZhy{}}\PYG{o}{\PYGZhy{}}\PYG{o}{\PYGZhy{}}\PYG{o}{\PYGZhy{}}\PYG{o}{\PYGZhy{}}\PYG{o}{\PYGZhy{}}\PYG{o}{\PYGZhy{}}\PYG{o}{\PYGZhy{}}\PYG{o}{\PYGZhy{}}\PYG{o}{\PYGZhy{}}\PYG{o}{\PYGZhy{}}\PYG{o}{\PYGZhy{}}\PYG{o}{\PYGZhy{}}\PYG{o}{\PYGZhy{}}\PYG{o}{\PYGZhy{}}\PYG{o}{\PYGZhy{}}\PYG{o}{+}
                                       \PYG{o}{\PYGZca{}}
                                       \PYG{o}{|} \PYG{n}{X2}\PYG{o}{/}\PYG{n}{Xn} \PYG{n}{Interface}
                                       \PYG{n}{v}
                     \PYG{o}{+}\PYG{o}{\PYGZhy{}}\PYG{o}{\PYGZhy{}}\PYG{o}{\PYGZhy{}}\PYG{o}{\PYGZhy{}}\PYG{o}{\PYGZhy{}}\PYG{o}{\PYGZhy{}}\PYG{o}{\PYGZhy{}}\PYG{o}{\PYGZhy{}}\PYG{o}{\PYGZhy{}}\PYG{o}{\PYGZhy{}}\PYG{o}{\PYGZhy{}}\PYG{o}{\PYGZhy{}}\PYG{o}{\PYGZhy{}}\PYG{o}{\PYGZhy{}}\PYG{o}{\PYGZhy{}}\PYG{o}{\PYGZhy{}}\PYG{o}{\PYGZhy{}}\PYG{o}{\PYGZhy{}}\PYG{o}{\PYGZhy{}}\PYG{o}{\PYGZhy{}}\PYG{o}{\PYGZhy{}}\PYG{o}{\PYGZhy{}}\PYG{o}{\PYGZhy{}}\PYG{o}{\PYGZhy{}}\PYG{o}{\PYGZhy{}}\PYG{o}{\PYGZhy{}}\PYG{o}{\PYGZhy{}}\PYG{o}{\PYGZhy{}}\PYG{o}{\PYGZhy{}}\PYG{o}{\PYGZhy{}}\PYG{o}{\PYGZhy{}}\PYG{o}{\PYGZhy{}}\PYG{o}{\PYGZhy{}}\PYG{o}{\PYGZhy{}}\PYG{o}{\PYGZhy{}}\PYG{o}{+}
                     \PYG{o}{|} \PYG{n}{Adjacent} \PYG{n}{Base} \PYG{n}{Stations} \PYG{p}{(}\PYG{l+m+mi}{10} \PYG{n}{Cells}\PYG{p}{)} \PYG{o}{|}
                     \PYG{o}{+}\PYG{o}{\PYGZhy{}}\PYG{o}{\PYGZhy{}}\PYG{o}{\PYGZhy{}}\PYG{o}{\PYGZhy{}}\PYG{o}{\PYGZhy{}}\PYG{o}{\PYGZhy{}}\PYG{o}{\PYGZhy{}}\PYG{o}{\PYGZhy{}}\PYG{o}{\PYGZhy{}}\PYG{o}{\PYGZhy{}}\PYG{o}{\PYGZhy{}}\PYG{o}{\PYGZhy{}}\PYG{o}{\PYGZhy{}}\PYG{o}{\PYGZhy{}}\PYG{o}{\PYGZhy{}}\PYG{o}{\PYGZhy{}}\PYG{o}{\PYGZhy{}}\PYG{o}{\PYGZhy{}}\PYG{o}{\PYGZhy{}}\PYG{o}{\PYGZhy{}}\PYG{o}{\PYGZhy{}}\PYG{o}{\PYGZhy{}}\PYG{o}{\PYGZhy{}}\PYG{o}{\PYGZhy{}}\PYG{o}{\PYGZhy{}}\PYG{o}{\PYGZhy{}}\PYG{o}{\PYGZhy{}}\PYG{o}{\PYGZhy{}}\PYG{o}{\PYGZhy{}}\PYG{o}{\PYGZhy{}}\PYG{o}{\PYGZhy{}}\PYG{o}{\PYGZhy{}}\PYG{o}{\PYGZhy{}}\PYG{o}{\PYGZhy{}}\PYG{o}{\PYGZhy{}}\PYG{o}{+}
\end{sphinxVerbatim}


\subsection{Infrastructure Specifications}
\label{\detokenize{deployement:infrastructure-specifications}}\begin{itemize}
\item {} 
\sphinxAtStartPar
\sphinxstylestrong{Host Platform:} Carrier\sphinxhyphen{}grade Edge Compute Node integrated within the gNodeB Centralized Unit Control Plane (CU\sphinxhyphen{}CP).

\item {} 
\sphinxAtStartPar
\sphinxstylestrong{Runtime Environment:} Docker Containerized Environment running an \sphinxstylestrong{ONNX Runtime} optimized for low\sphinxhyphen{}latency inference (\textless{} 2ms per execution loop).

\item {} 
\sphinxAtStartPar
\sphinxstylestrong{Input Vector Data:} Real\sphinxhyphen{}time metrics streaming from the cell Layer 3 filtering:
\begin{itemize}
\item {} 
\sphinxAtStartPar
Reference Signal Received Power (RSRP) histories for all 10 candidate cells.

\item {} 
\sphinxAtStartPar
Reference Signal Received Quality (RSRQ) histories.

\item {} 
\sphinxAtStartPar
Channel Quality Indicators (CQI).

\item {} 
\sphinxAtStartPar
Drone load/capacity telemetry.

\end{itemize}

\end{itemize}


\bigskip\hrule\bigskip



\section{5. Multi\sphinxhyphen{}Horizon Operational Pipeline (250ms Execution Loop)}
\label{\detokenize{deployement:multi-horizon-operational-pipeline-250ms-execution-loop}}
\sphinxAtStartPar
The system enforces a strict real\sphinxhyphen{}time telemetry\sphinxhyphen{}to\sphinxhyphen{}execution pipeline designed to feed the 3GPP standard signaling loop.

\begin{sphinxVerbatim}[commandchars=\\\{\}]
       \PYG{p}{[}\PYG{n}{T} \PYG{o}{=} \PYG{l+m+mi}{0}\PYG{n}{ms}\PYG{p}{]} \PYG{n}{Stream} \PYG{n}{L3} \PYG{n}{filtered} \PYG{n}{metrics} \PYG{p}{(}\PYG{n}{RSRP}\PYG{o}{/}\PYG{n}{RSRQ}\PYG{o}{/}\PYG{n}{CQI}\PYG{p}{)} \PYG{k+kn}{from}\PYG{+w}{ }\PYG{l+m+mi}{10} \PYG{n}{cells}
                           \PYG{o}{|}
                           \PYG{n}{v}
       \PYG{p}{[}\PYG{n}{T} \PYG{o}{+} \PYG{l+m+mi}{2}\PYG{n}{ms}\PYG{p}{]} \PYG{n}{Vectorization} \PYG{o}{\PYGZam{}} \PYG{n}{Normalization} \PYG{n}{into} \PYG{n}{Time}\PYG{o}{\PYGZhy{}}\PYG{n}{Series} \PYG{n}{Matrix}
                           \PYG{o}{|}
                           \PYG{n}{v}
       \PYG{p}{[}\PYG{n}{T} \PYG{o}{+} \PYG{l+m+mi}{4}\PYG{n}{ms}\PYG{p}{]} \PYG{n}{Execute} \PYG{n}{Inference} \PYG{n}{Engine} \PYG{p}{(}\PYG{n}{Multi}\PYG{o}{\PYGZhy{}}\PYG{n}{Horizon} \PYG{l+m+mi}{5}\PYG{o}{\PYGZhy{}}\PYG{n}{step} \PYG{n}{Prediction}\PYG{p}{)}
                           \PYG{o}{|}
                           \PYG{n}{v}
       \PYG{p}{[}\PYG{n}{T} \PYG{o}{+} \PYG{l+m+mi}{5}\PYG{n}{ms}\PYG{p}{]} \PYG{n}{Extract} \PYG{n}{Top}\PYG{o}{\PYGZhy{}}\PYG{l+m+mi}{3} \PYG{n}{Predicted} \PYG{n}{Candidate} \PYG{n}{Cells} \PYG{p}{(}\PYG{l+m+mi}{85}\PYG{o}{\PYGZpc{}} \PYG{n}{Accuracy}\PYG{p}{)}
                           \PYG{o}{|}
                           \PYG{n}{v}
       \PYG{p}{[}\PYG{n}{T} \PYG{o}{+} \PYG{l+m+mi}{7}\PYG{n}{ms}\PYG{p}{]} \PYG{n}{gNodeB} \PYG{n}{issues} \PYG{n}{RRC} \PYG{n}{Reconfiguration} \PYG{k}{with} \PYG{n}{CHO} \PYG{n}{Configurations}
                           \PYG{o}{|}
                           \PYG{n}{v}
       \PYG{p}{[}\PYG{n}{T} \PYG{o}{+} \PYG{l+m+mi}{250}\PYG{n}{ms}\PYG{p}{]} \PYG{n}{UE} \PYG{n}{encounters} \PYG{n}{real}\PYG{o}{\PYGZhy{}}\PYG{n}{time} \PYG{n}{channel} \PYG{n}{change} \PYG{o}{\PYGZhy{}}\PYG{o}{\PYGZgt{}} \PYG{n}{Instantly} \PYG{n}{connects} 
                   \PYG{n}{to} \PYG{n}{pre}\PYG{o}{\PYGZhy{}}\PYG{n}{prepared} \PYG{n}{target} \PYG{n}{cell} \PYG{n}{without} \PYG{n}{signaling} \PYG{n}{delay}
\end{sphinxVerbatim}


\subsection{Detailed Pipeline Stages}
\label{\detokenize{deployement:detailed-pipeline-stages}}\begin{enumerate}
\sphinxsetlistlabels{\arabic}{enumi}{enumii}{}{.}%
\item {} 
\sphinxAtStartPar
\sphinxstylestrong{Data Ingestion (Every 50ms):} The CU\sphinxhyphen{}CP extracts the sliding window of the last 10 historical timestamps (500ms total history) of RSRP/RSRQ measurements for the target UE across the 10 cells.

\item {} 
\sphinxAtStartPar
\sphinxstylestrong{Inference Execution:} The time\sphinxhyphen{}series matrix is fed into the model. The model computes the multi\sphinxhyphen{}horizon probability distribution for the cell states at \(T+50	ext{ms}, T+100	ext{ms}, T+150	ext{ms}, T+200	ext{ms}, 	ext{ and } T+250	ext{ms}\).

\item {} 
\sphinxAtStartPar
\sphinxstylestrong{Top\sphinxhyphen{}3 Target Selection:} The final output layer ranks the top 3 target cells projected at \(T+250	ext{ms}\).

\item {} 
\sphinxAtStartPar
\sphinxstylestrong{Proactive Resource Preparation:} If the predicted top\sphinxhyphen{}1 cell differs from the serving cell and crosses the safety delta threshold, the gNodeB triggers standard X2/Xn interface preparation signaling to the top 3 target neighbors. Context vectors and target RACH pre\sphinxhyphen{}allocations are completed \sphinxstyleemphasis{before} the 250ms window expires.

\end{enumerate}


\bigskip\hrule\bigskip



\section{6. Deployment Verification \& Fallback Controls}
\label{\detokenize{deployement:deployment-verification-fallback-controls}}
\sphinxAtStartPar
Because this engine runs within a critical telecommunications production cluster, strict guardrails must ensure service continuity if unexpected field anomalies occur.


\subsection{KPI Guardrail Triggers}
\label{\detokenize{deployement:kpi-guardrail-triggers}}\begin{itemize}
\item {} 
\sphinxAtStartPar
\sphinxstylestrong{Anomaly Detector:} If the current real\sphinxhyphen{}time Top\sphinxhyphen{}1 measurement mismatches the 250ms historical prediction for 3 consecutive execution loops, a local prediction confidence degradation flag is raised.

\item {} 
\sphinxAtStartPar
\sphinxstylestrong{Fallback Protocol:} Upon flag activation, the gNodeB temporarily falls back to standard \sphinxstylestrong{3GPP Release 15 Event A3 (Neighbor becomes offset better than spcell)} reactive handover logic. This bypasses the predictive engine until the context alignment score stabilizes above 80\%.

\item {} 
\sphinxAtStartPar
\sphinxstylestrong{Drone Load Balancing:} If a drone node’s load exceeds 85\% capacity, the predictive inference engine automatically appends a penalty scalar to that cell’s predicted score, naturally shedding user handovers to adjacent ground macro cells.

\end{itemize}


\bigskip\hrule\bigskip



\section{7. Next\sphinxhyphen{}Step Integration Milestone Checklist}
\label{\detokenize{deployement:next-step-integration-milestone-checklist}}\begin{itemize}
\item {} 
\sphinxAtStartPar
{[}{]} \sphinxstylestrong{Phase 1:} Export fine\sphinxhyphen{}tuned weights from Python/Sionna framework to \sphinxcode{\sphinxupquote{.onnx}} format.

\item {} 
\sphinxAtStartPar
{[} {]} \sphinxstylestrong{Phase 2:} Integrate ONNX model file into the testbench gNodeB emulator.

\item {} 
\sphinxAtStartPar
{[} {]} \sphinxstylestrong{Phase 3:} Establish live telemetry pipelines using standard RAN Intelligent Controller (RIC) or direct southbound interfaces to pull RSRP streams.

\item {} 
\sphinxAtStartPar
{[} {]} \sphinxstylestrong{Phase 4:} Conduct shadow\sphinxhyphen{}mode testing on production NR servers (predicting without executing signaling changes) to validate the 85\% Top\sphinxhyphen{}3 accuracy against real\sphinxhyphen{}world user equipment.

\item {} 
\sphinxAtStartPar
{[} {]} \sphinxstylestrong{Phase 5:} Activate closed\sphinxhyphen{}loop execution for selected User Equipment groups to achieve seamless predictive handovers.

\end{itemize}



\renewcommand{\indexname}{Index}
\printindex
\end{document}